% Options for packages loaded elsewhere
% Options for packages loaded elsewhere
\PassOptionsToPackage{unicode}{hyperref}
\PassOptionsToPackage{hyphens}{url}
\PassOptionsToPackage{dvipsnames,svgnames,x11names}{xcolor}
%
\documentclass[
  dutch,
  letterpaper,
  DIV=11,
  numbers=noendperiod]{scrartcl}
\usepackage{xcolor}
\usepackage{amsmath,amssymb}
\setcounter{secnumdepth}{-\maxdimen} % remove section numbering
\usepackage{iftex}
\ifPDFTeX
  \usepackage[T1]{fontenc}
  \usepackage[utf8]{inputenc}
  \usepackage{textcomp} % provide euro and other symbols
\else % if luatex or xetex
  \usepackage{unicode-math} % this also loads fontspec
  \defaultfontfeatures{Scale=MatchLowercase}
  \defaultfontfeatures[\rmfamily]{Ligatures=TeX,Scale=1}
\fi
\usepackage{lmodern}
\ifPDFTeX\else
  % xetex/luatex font selection
\fi
% Use upquote if available, for straight quotes in verbatim environments
\IfFileExists{upquote.sty}{\usepackage{upquote}}{}
\IfFileExists{microtype.sty}{% use microtype if available
  \usepackage[]{microtype}
  \UseMicrotypeSet[protrusion]{basicmath} % disable protrusion for tt fonts
}{}
\makeatletter
\@ifundefined{KOMAClassName}{% if non-KOMA class
  \IfFileExists{parskip.sty}{%
    \usepackage{parskip}
  }{% else
    \setlength{\parindent}{0pt}
    \setlength{\parskip}{6pt plus 2pt minus 1pt}}
}{% if KOMA class
  \KOMAoptions{parskip=half}}
\makeatother
% Make \paragraph and \subparagraph free-standing
\makeatletter
\ifx\paragraph\undefined\else
  \let\oldparagraph\paragraph
  \renewcommand{\paragraph}{
    \@ifstar
      \xxxParagraphStar
      \xxxParagraphNoStar
  }
  \newcommand{\xxxParagraphStar}[1]{\oldparagraph*{#1}\mbox{}}
  \newcommand{\xxxParagraphNoStar}[1]{\oldparagraph{#1}\mbox{}}
\fi
\ifx\subparagraph\undefined\else
  \let\oldsubparagraph\subparagraph
  \renewcommand{\subparagraph}{
    \@ifstar
      \xxxSubParagraphStar
      \xxxSubParagraphNoStar
  }
  \newcommand{\xxxSubParagraphStar}[1]{\oldsubparagraph*{#1}\mbox{}}
  \newcommand{\xxxSubParagraphNoStar}[1]{\oldsubparagraph{#1}\mbox{}}
\fi
\makeatother


\usepackage{longtable,booktabs,array}
\usepackage{calc} % for calculating minipage widths
% Correct order of tables after \paragraph or \subparagraph
\usepackage{etoolbox}
\makeatletter
\patchcmd\longtable{\par}{\if@noskipsec\mbox{}\fi\par}{}{}
\makeatother
% Allow footnotes in longtable head/foot
\IfFileExists{footnotehyper.sty}{\usepackage{footnotehyper}}{\usepackage{footnote}}
\makesavenoteenv{longtable}
\usepackage{graphicx}
\makeatletter
\newsavebox\pandoc@box
\newcommand*\pandocbounded[1]{% scales image to fit in text height/width
  \sbox\pandoc@box{#1}%
  \Gscale@div\@tempa{\textheight}{\dimexpr\ht\pandoc@box+\dp\pandoc@box\relax}%
  \Gscale@div\@tempb{\linewidth}{\wd\pandoc@box}%
  \ifdim\@tempb\p@<\@tempa\p@\let\@tempa\@tempb\fi% select the smaller of both
  \ifdim\@tempa\p@<\p@\scalebox{\@tempa}{\usebox\pandoc@box}%
  \else\usebox{\pandoc@box}%
  \fi%
}
% Set default figure placement to htbp
\def\fps@figure{htbp}
\makeatother


% definitions for citeproc citations
\NewDocumentCommand\citeproctext{}{}
\NewDocumentCommand\citeproc{mm}{%
  \begingroup\def\citeproctext{#2}\cite{#1}\endgroup}
\makeatletter
 % allow citations to break across lines
 \let\@cite@ofmt\@firstofone
 % avoid brackets around text for \cite:
 \def\@biblabel#1{}
 \def\@cite#1#2{{#1\if@tempswa , #2\fi}}
\makeatother
\newlength{\cslhangindent}
\setlength{\cslhangindent}{1.5em}
\newlength{\csllabelwidth}
\setlength{\csllabelwidth}{3em}
\newenvironment{CSLReferences}[2] % #1 hanging-indent, #2 entry-spacing
 {\begin{list}{}{%
  \setlength{\itemindent}{0pt}
  \setlength{\leftmargin}{0pt}
  \setlength{\parsep}{0pt}
  % turn on hanging indent if param 1 is 1
  \ifodd #1
   \setlength{\leftmargin}{\cslhangindent}
   \setlength{\itemindent}{-1\cslhangindent}
  \fi
  % set entry spacing
  \setlength{\itemsep}{#2\baselineskip}}}
 {\end{list}}
\usepackage{calc}
\newcommand{\CSLBlock}[1]{\hfill\break\parbox[t]{\linewidth}{\strut\ignorespaces#1\strut}}
\newcommand{\CSLLeftMargin}[1]{\parbox[t]{\csllabelwidth}{\strut#1\strut}}
\newcommand{\CSLRightInline}[1]{\parbox[t]{\linewidth - \csllabelwidth}{\strut#1\strut}}
\newcommand{\CSLIndent}[1]{\hspace{\cslhangindent}#1}

\ifLuaTeX
\usepackage[bidi=basic]{babel}
\else
\usepackage[bidi=default]{babel}
\fi
% get rid of language-specific shorthands (see #6817):
\let\LanguageShortHands\languageshorthands
\def\languageshorthands#1{}


\setlength{\emergencystretch}{3em} % prevent overfull lines

\providecommand{\tightlist}{%
  \setlength{\itemsep}{0pt}\setlength{\parskip}{0pt}}



 


\KOMAoption{captions}{tableheading}
\makeatletter
\@ifpackageloaded{caption}{}{\usepackage{caption}}
\AtBeginDocument{%
\ifdefined\contentsname
  \renewcommand*\contentsname{Inhoudsopgave}
\else
  \newcommand\contentsname{Inhoudsopgave}
\fi
\ifdefined\listfigurename
  \renewcommand*\listfigurename{Lijst van figuren}
\else
  \newcommand\listfigurename{Lijst van figuren}
\fi
\ifdefined\listtablename
  \renewcommand*\listtablename{Lijst van tabellen}
\else
  \newcommand\listtablename{Lijst van tabellen}
\fi
\ifdefined\figurename
  \renewcommand*\figurename{Figuur}
\else
  \newcommand\figurename{Figuur}
\fi
\ifdefined\tablename
  \renewcommand*\tablename{Tabel}
\else
  \newcommand\tablename{Tabel}
\fi
}
\@ifpackageloaded{float}{}{\usepackage{float}}
\floatstyle{ruled}
\@ifundefined{c@chapter}{\newfloat{codelisting}{h}{lop}}{\newfloat{codelisting}{h}{lop}[chapter]}
\floatname{codelisting}{Listing}
\newcommand*\listoflistings{\listof{codelisting}{Lijst van listings}}
\makeatother
\makeatletter
\makeatother
\makeatletter
\@ifpackageloaded{caption}{}{\usepackage{caption}}
\@ifpackageloaded{subcaption}{}{\usepackage{subcaption}}
\makeatother
\usepackage{bookmark}
\IfFileExists{xurl.sty}{\usepackage{xurl}}{} % add URL line breaks if available
\urlstyle{same}
\hypersetup{
  pdftitle={DigComp 3.0 Nederlands},
  pdfauthor={Pierre Gorissen, Manon van Zanten},
  pdflang={nl},
  colorlinks=true,
  linkcolor={blue},
  filecolor={Maroon},
  citecolor={Blue},
  urlcolor={Blue},
  pdfcreator={LaTeX via pandoc}}


\title{DigComp 3.0 Nederlands}
\author{Pierre Gorissen, Manon van Zanten}
\date{2026-03-06}
\begin{document}
\maketitle


\subsection*{Colofon}\label{sec-colofon}
\addcontentsline{toc}{subsection}{Colofon}

\textbf{Originele bron:} Cosgrove, J. and Cachia, R., DigComp 3.0:
European Digital Competence Framework - Fifth Edition, Publications
Office of the European Union, Luxembourg, 2025,
https://data.europa.eu/doi/10.2760/0001149, JRC144121.

De Europese Commissie is niet verantwoordelijk voor deze vertaalde
versie en de bijbehorende databestanden die in het Nederlands
beschikbaar zijn. De Europese Commissie staat de vertaling van DigComp
3.0 in andere talen toe, maar onderschrijft de gewijzigde, aangepaste of
vertaalde versies niet.

Het hergebruikbeleid van de documenten van de Europese Commissie wordt
beschreven in Besluit 2011/833/EU. Tenzij anders vermeld, is hergebruik
van het originele DigComp-document toegestaan onder de CC BY 4.0
licentie.

De vertaling naar het Nederlands is uitgevoerd onder auspiciën van het
iXperium Centre of Expertise Teaching and Learning. Er zijn geen
wijzigingen aangebracht in de originele competentiegebieden, behalve een
vertaling waarbij de intentie zo dicht mogelijk bij het origineel is
gebleven.

Maak bij gebruik van deze vertaling vermelding van de volgende
referentie: Gorissen, P. \& van Zanten, M. (2026). \emph{DigComp 3.3:
Europees kader voor digitale competenties}. iXperium Centre of Expertise
Leren met ict.

\textbf{\texttt{{[}standaard\ CC-BY\ logo\ iXperium\ invoegen{]}}}

\subsection*{Over deze vertaling}\label{sec-over-vertaling}
\addcontentsline{toc}{subsection}{Over deze vertaling}

Deze Nederlandse vertaling beperkt zich tot de belangrijkste onderdelen
van de DigComp 3.0. De vertaling bevat een Inleiding op basis van de
Engelse tekst, het DigComp 3.0 raamwerk, de competentiegebieden en
competenties, beheersingsniveaus en leeruitkomsten.

\newpage{}

\section{Inleiding}\label{sec-inleiding}

\subsection{Over DigComp 3.0}\label{sec-over-digcomp-3.0}

DigComp 3.0 is de vijfde editie van het Europese Digitale
competentie-raamwerk. Het beschrijft kennis, vaardigheden en houdingen
die nodig zijn om digitaal bekwaam te zijn voor het dagelijks leven,
deelname aan de samenleving, werken en leren. Het raamwerk is
technologie-neutraal en ontworpen om op maat gemaakt te worden voor
onderwijs- en arbeidscontexten.

Het ondersteunt EU-beleidslijnen zoals de Unie van Vaardigheden en het
Beleidsprogramma Digitale Decennium. DigComp 3.0 bevat ontwikkelingen
die sinds 2022 hebben plaatsgevonden, inclusief de integratie van
AI-competenties.

\subsection{Het belang van Digitale
Competentie}\label{het-belang-van-digitale-competentie}

Digitale omgevingen zijn verweven met het dagelijks leven, wat zowel
voordelen als risico's met zich meebrengt. Het is essentieel om tekorten
in vaardigheden te overbruggen.

Digitale competenties bieden individuen voordelen zoals het gebruik van
publieke platforms, privacybeheer en welzijn (Morte-Nadal \&
Esteban-Navarro, 2025; Pouliakas et al., 2025; Stalmach et al., 2023;
Țarcă et al., 2024), en vormen een middel om sociale uitsluiting aan te
pakken (Boerkamp et al., 2024; Brundle et al., 2025) en de
concurrentiekracht te versterken (Draghi, 2024; Pakhnenko et al., 2025).

In beleid en initiatieven worden `digitale vaardigheden' en `digitale
competentie' vaak door elkaar gebruikt. In DigComp 3.0 wordt de term
`digitale competentie' gebruikt om kennis, vaardigheden en houdingen te
omvatten.

\subsection{Digitale feiten en
cijfers}\label{digitale-feiten-en-cijfers}

\begin{longtable}[]{@{}
  >{\raggedright\arraybackslash}p{(\linewidth - 4\tabcolsep) * \real{0.3333}}
  >{\raggedright\arraybackslash}p{(\linewidth - 4\tabcolsep) * \real{0.3333}}
  >{\raggedright\arraybackslash}p{(\linewidth - 4\tabcolsep) * \real{0.3333}}@{}}
\caption{Digitale competentie -- een selectie van feiten en
cijfers.}\label{tbl-stats}\tabularnewline
\toprule\noalign{}
\begin{minipage}[b]{\linewidth}\raggedright
Kinderen en jongeren
\end{minipage} & \begin{minipage}[b]{\linewidth}\raggedright
Volwassenen in het algemeen
\end{minipage} & \begin{minipage}[b]{\linewidth}\raggedright
Werkenden
\end{minipage} \\
\midrule\noalign{}
\endfirsthead
\toprule\noalign{}
\begin{minipage}[b]{\linewidth}\raggedright
Kinderen en jongeren
\end{minipage} & \begin{minipage}[b]{\linewidth}\raggedright
Volwassenen in het algemeen
\end{minipage} & \begin{minipage}[b]{\linewidth}\raggedright
Werkenden
\end{minipage} \\
\midrule\noalign{}
\endhead
\bottomrule\noalign{}
\endlastfoot
43\% van de leerlingen bereikte in 2023 geen basisniveau (Duckworth \&
Fraillon, 2025). & 56\% van de volwassenen beschikte in 2023 over
digitale basisvaardigheden (European Commission, 2025a). & 92\% van de
werknemers gebruikte in 2024-2025 digitale apparaten (Eurostat,
2025). \\
96\% van de 15-jarigen gebruikte in 2022 dagelijks sociale media
(European Commission, 2025b). & 70\% van de volwassenen had in 2024
digitale interactie met de overheid (European Commission, 2025a). & 30\%
van de werknemers gebruikte AI-systemen op het werk in 2024-2025
(Pouliakas et al., 2025). \\
14\%-16\% was slachtoffer van cyberpesten in 2022 (European Commission,
2025b). & 49\% meldde in 2023 onjuiste inhoud op sociale media (European
Commission, 2025a). & 35\% moest in 2020-2021 nieuwe technologie leren
voor werk (European Commission, 2023a). \\
9-12\% rapporteerde problematisch gebruik van sociale media (European
Commission, 2025b). & 58\% zocht in 2022 naar gezondheidsinformatie
(European Commission, 2025a). & 42\% rapporteerde in 2024 een kloof in
AI-vaardigheden (Pouliakas et al., 2025). \\
\end{longtable}

Tabel~\ref{tbl-stats} illustreert enkele implicaties van digitale
technologie voor digitale competentie bij kinderen en jongeren,
volwassenen in het algemeen en werknemers, aan de hand van een selectie
van statistisch bewijs uit recente enquêtes.

\subsection{Adoptie en gebruik van
DigComp}\label{sec-adoptie-en-gebruik-van-digcomp}

Een belangrijk kenmerk van DigComp is de mate waarin experts en
belanghebbenden uit Europa en daarbuiten bijdragen aan de ontwikkeling
ervan. Dit helpt de kwaliteit en relevantie te waarborgen en ondersteunt
de adoptie ervan. DigComp wordt veelvuldig gebruikt in Europa en
daarbuiten in contexten van werkgelegenheid, onderwijs en opleiding
(Centeno et al., 2024b; Centeno \& Cosgrove, 2025; Kluzer et al., 2020),
en vormt de conceptuele basis van de Digital Skills Indicator (DSI)
(Vuorikari et al., 2022), die wordt gebruikt om het niveau van digitale
basisvaardigheden in Europa te monitoren via het beleidsprogramma voor
het digitale decennium. Naast het bieden van een gemeenschappelijk
begrip van digitale competentie, wordt DigComp onder andere gebruikt om
Europees, nationaal en regionaal beleid te sturen; beoordelingen te
ontwikkelen; de transparantie of vergelijkbaarheid van onderwijs- en
opleidingscursussen te verbeteren; leerprestaties te erkennen of te
valideren (bijvoorbeeld door certificering van digitale vaardigheden);
en profielen van digitale competentie in specifieke banen of rollen te
definiëren.

\subsection{Waarden en principes die richting geven aan DigComp
3.0}\label{sec-waarden-en-principes}

De visie voor DigComp 3.0 is dat het een verenigend, coherent, helder,
relevant en actueel beeld geeft van digitale competentie dat voortbouwt
op eerdere versies, en dat kan worden gebruikt door een verscheidenheid
aan belanghebbenden die het gemeenschappelijke doel delen om behoeften
op het gebied van digitale competentie te begrijpen en te herkennen en
de ontwikkeling ervan te ondersteunen.

DigComp 3.0, de vijfde editie van het kader, blijft het belangrijkste
referentiekader voor digitale competentie in de EU. Digitale competentie
is een van de acht sleutelcompetenties van de aanbeveling van de Raad
inzake sleutelcompetenties voor een leven lang leren (European
Commission, 2018).

\emph{Kerncompetenties zijn de kennis, vaardigheden en houdingen die
iedereen nodig heeft voor persoonlijke ontplooiing en ontwikkeling,
inzetbaarheid op de arbeidsmarkt, sociale inclusie en actief
burgerschap.}

DigComp 3.0 belichaamt de waarden van de Europese verklaring over
digitale rechten en beginselen voor het digitale decennium (European
Commission, 2023b), die een mensgerichte visie op de digitale
transformatie bevordert. Het is georganiseerd in zes thema's (zie figuur
2). De verklaring bouwt voort op het EU-Handvest van de grondrechten
(European Commission, 2012).

De zes thema's van de Europese verklaring over digitale rechten en
beginselen (European Commission, 2023b).

\begin{itemize}
\item
  \textbf{Mensen centraal:} Digitale technologie moeten de rechten van
  mensen beschermen, de democratie ondersteunen en ervoor zorgen dat
  alle digitale actoren verantwoordelijk en veilig handelen. De EU
  promoot deze waarden over de hele wereld.
\item
  \textbf{Solidariteit en inclusie:} Technologie moet mensen verenigen,
  niet verdelen. Iedereen moet toegang hebben tot internet, tot digitale
  vaardigheden, tot digitale overheidsdiensten en tot eerlijke
  arbeidsomstandigheden.
\item
  \textbf{Vrijheid van keuze:} Mensen moeten kunnen profiteren van een
  eerlijke online omgeving, veilig zijn voor illegale en schadelijke
  inhoud, en in hun kracht worden gezet wanneer ze interageren met
  nieuwe en evoluerende technologie zoals kunstmatige intelligentie.
\item
  \textbf{Participatie:} Burgers moeten kunnen deelnemen aan het
  democratische proces op alle niveaus en controle hebben over hun eigen
  gegevens.
\item
  \textbf{Veiligheid en beveiliging:} De digitale omgeving moet veilig
  en betrouwbaar zijn. Alle gebruikers, van kind tot op hoge leeftijd,
  moeten worden gesterkt en beschermd.
\item
  \textbf{Duurzaamheid:} Digitale apparaten moeten duurzaamheid en de
  groene transitie ondersteunen. Mensen moeten weten wat de
  milieu-impact en het energieverbruik van hun apparaten zijn.
\end{itemize}

\emph{Bron: Gebaseerd op https://digital-strategy.ec.europa.eu/en/fact-
pages/digital-rights-and-principles}

\subsection{Wat is er nieuw in DigComp
3.0?}\label{sec-wat-is-er-nieuw-in-digcomp-3.0}

Verschillende belangrijke thema's met betrekking tot de inhoud (d.w.z.
waaruit digitale competentie bestaat) en toepassing (d.w.z. hoe het
raamwerk wordt aangepast en gebruikt en de rol ervan in onderwijs,
opleiding en werkgelegenheidssystemen) gaven richting aan de
ontwikkeling van DigComp 3.0 (Tabel~\ref{tbl-themas}). Deze werden
geïdentificeerd op basis van beleids- en academische literatuur en
overleg met experts en belanghebbenden en zijn holistisch ingebed in de
inhoud en het ontwerp van het raamwerk.

\begin{longtable}[]{@{}
  >{\raggedright\arraybackslash}p{(\linewidth - 0\tabcolsep) * \real{1.0000}}@{}}
\caption{Inhouds- en toepassings- thema's die richting geven aan DigComp
3.0 ontwikkeling.}\label{tbl-themas}\tabularnewline
\toprule\noalign{}
\endfirsthead
\endhead
\bottomrule\noalign{}
\endlastfoot
\textbf{Inhoudsthema's} \\
Kunstmatige intelligentie (inclusief generatieve AI) competentie \\
Cyberbeveiliging competentie \\
Digitale rechten, keuze en verantwoordelijkheden \\
Welzijn in digitale omgevingen \\
Competentie om misinformatie en desinformatie aan te pakken \\
\textbf{Toepassingsthema's} \\
Digitale competentie als essentieel onderdeel van een leven lang
leren \\
Erkenning van randvoorwaarden voor het verwerven van digitale
competentie op basisniveau \\
Erkenning van verschillen in behoeften aan digitale competentie tussen
personen en door de tijd heen \\
Behoefte aan flexibele, wendbare toepassingen van het kader \\
\end{longtable}

\emph{Bron: Eigen bewerking door JRC.}

Tabel~\ref{tbl-doelen} vat de doelstellingen en belangrijkste kenmerken
van de DigComp 3.0-update samen. De update van DigComp wordt gedreven
door aanzienlijke technologische ontwikkelingen, trends en praktijken
(zoals de snelle verspreiding van generatieve AI) die hebben
plaatsgevonden sinds de publicatie van DigComp 2.2 in 2022. Deze hebben
verreikende implicaties voor de digitale competenties van individuen en
komen tot uiting in beleidsprioriteiten en zorgen van belanghebbenden
(Abendroth-Dias et al., 2025; Farias-Gaytan et al., 2023; Lewandowsky et
al., 2020; Onesi-Ozigagun et al., 2024). DigComp 3.0 reageert ook op
verzoeken van bestaande gebruikers en belanghebbenden, die meer
duidelijkheid wensten over praktische toepassingen van DigComp. Een
belangrijke aanbeveling in een studie naar de haalbaarheid van een
Europees certificaat voor digitale vaardigheden, waarbij 650
belanghebbenden uit alle EU-lidstaten betrokken waren, was om
leerresultaten voor DigComp te ontwikkelen (Centeno et al., 2024a).

\begin{longtable}[]{@{}
  >{\raggedright\arraybackslash}p{(\linewidth - 2\tabcolsep) * \real{0.5000}}
  >{\raggedright\arraybackslash}p{(\linewidth - 2\tabcolsep) * \real{0.5000}}@{}}
\caption{DigComp 3.0: doelstellingen, ontwikkelingen en middelen voor
implementatie.}\label{tbl-doelen}\tabularnewline
\toprule\noalign{}
\begin{minipage}[b]{\linewidth}\raggedright
Categorie
\end{minipage} & \begin{minipage}[b]{\linewidth}\raggedright
Details
\end{minipage} \\
\midrule\noalign{}
\endfirsthead
\toprule\noalign{}
\begin{minipage}[b]{\linewidth}\raggedright
Categorie
\end{minipage} & \begin{minipage}[b]{\linewidth}\raggedright
Details
\end{minipage} \\
\midrule\noalign{}
\endhead
\bottomrule\noalign{}
\endlastfoot
\textbf{DOELSTELLINGEN} & 1. Integreer recente en opkomende digitale
technologische trends en hun implicaties voor digitale competentie,
terwijl de algemene structuur van het raamwerk en technologische
neutraliteit behouden blijven. \\
& 2. Ontwikkeling van leerresultaten en andere passende verbeteringen
ter ondersteuning van duidelijkheid en operationele afstemming in de
toepassing van DigComp. \\
\textbf{WIJZIGINGEN EN ONTWIKKELINGEN} & • Updates aan de bewoordingen
van de vijf competentiegebieden en 21 competenties om de huidige
technologische trends te weerspiegelen. \\
& • Een nieuwe algemene beschrijving van beheersingsniveaus (Basis,
Gevorderd, Geavanceerd en Zeer geavanceerd) die aan de vorige versie
kunnen worden gekoppeld, en nieuwe en herziene competentiebeschrijvingen
voor elk beheersingsniveau en competentie. \\
& • Ontwikkeling van leerresultaten, voor een duidelijkere interpretatie
en implementatie. \\
& • Een systematische, transversale integratie van AI-competentie in het
raamwerk die voortbouwt op DigComp 2.2, en daarnaast recente
ontwikkelingen met betrekking tot AI omvat. \\
\textbf{HULPMIDDELEN VOOR IMPLEMENTATIE} & • Een gedetailleerde
woordenlijst met circa 120 kerntermen in dit document. \\
& • Duidelijke en gebruiksvriendelijke informatie over DigComp 3.0 op de
JRC-DigComp-webruimte. \\
& • Ook op de JRC-DigComp-webruimte zijn versies van DigComp 3.0
beschikbaar in spreadsheet- en linked open data (JSON) indelingen. \\
\end{longtable}

\emph{Bron: eigen bewerking JRC.}

\subsection{DigComp 3.0 gebruiken}\label{sec-digcomp-3.0-gebruiken}

DigComp 3.0 is niet-voorschrijvend en is bedoeld om te worden aangepast
aan specifieke toepassingen en behoeften. DigComp is ontworpen om
initiatieven en acties te sturen en te inspireren die de ontwikkeling
van digitale competentie van individuen ondersteunen, zowel in het
algemeen (zoals in algemene beleidsvorming) als bij specifieke groepen
(zoals jonge lerenden, volwassen lerenden, kwetsbare of
gemarginaliseerde individuen, werknemers of werkzoekenden). Als zodanig
moet het worden beschouwd als een startpunt, dat moet worden aangepast
aan een specifiek doel.

\newpage{}

\section{DigComp 3.0
raamwerkcomponenten}\label{sec-digcomp-3.0-raamwerkcomponenten}

\subsection{Overzicht}\label{sec-overzicht}

Tabel~\ref{tbl-componenten} beschrijft de componenten van DigComp 3.0 en
hun doeleinden: definitie van digitale competentie, competentiegebieden,
competenties, niveaus van bekwaamheid en leerresultaten.

\begin{longtable}[]{@{}
  >{\raggedright\arraybackslash}p{(\linewidth - 2\tabcolsep) * \real{0.2500}}
  >{\raggedright\arraybackslash}p{(\linewidth - 2\tabcolsep) * \real{0.7500}}@{}}
\caption{Componenten van het DigComp 3.0
raamwerk}\label{tbl-componenten}\tabularnewline
\toprule\noalign{}
\begin{minipage}[b]{\linewidth}\raggedright
Component
\end{minipage} & \begin{minipage}[b]{\linewidth}\raggedright
Beschrijving en doel
\end{minipage} \\
\midrule\noalign{}
\endfirsthead
\toprule\noalign{}
\begin{minipage}[b]{\linewidth}\raggedright
Component
\end{minipage} & \begin{minipage}[b]{\linewidth}\raggedright
Beschrijving en doel
\end{minipage} \\
\midrule\noalign{}
\endhead
\bottomrule\noalign{}
\endlastfoot
\textbf{Definitie van digitale competentie} & Stelt de inhoud en de
reikwijdte van het kader vast, maakt onderscheid tussen kennis,
vaardigheden en houdingsaspecten van digitale competentie, en plaatst
het kader in de bredere context van sleutelcompetenties voor een leven
lang leren. Zie
Paragraaf~\ref{sec-definitie-van-digitale-competentie}. \\
\textbf{Competentiegebieden} & Organiseert competenties in vijf
thematische groepen, bestaande uit namen (titels) van
competentiegebieden en descriptoren (korte beschrijvingen van de
competenties die in elk competentiegebied vallen). Zie
Paragraaf~\ref{sec-competentiegebieden-en-competenties}. \\
\textbf{Competenties} & Definieert 21 competenties, georganiseerd onder
de vijf competentiegebieden, bestaande uit competentienamen (titels) en
descriptoren (korte beschrijvingen van waaruit elke competentie
bestaat). Zie
Paragraaf~\ref{sec-competentiegebieden-en-competenties}. \\
\textbf{Beheersingsniveaus} & Stelt een continuüm van beheersing vast op
basis van cognitieve belasting, taakcomplexiteit en mate van autonomie.
Algemene beheersingsniveaus (zie Paragraaf~\ref{sec-beheersingsniveaus})
beschrijven digitale competentie op de niveaus Basis, Gemiddeld,
Gevorderd en Zeer gevorderd. Competentieverklaringen (zie
\textbf{?@sec-competentieverklaringen}) bieden een gedetailleerder beeld
van de beheersingsniveaus voor elk van de 21 competenties. \\
\textbf{Leerresultaten} & Leerresultaten (zie
Paragraaf~\ref{sec-leerresultaten}) bieden het meest gedetailleerde
(`granulaire') beeld van het raamwerk en bestaan uit verklaringen over
wat een individu geacht wordt te weten of te kunnen na voltooiing van
het leerproces (van welke aard dan ook). Elk van de leerresultaten is
geclassificeerd naar competentie, bekwaamheidsniveau en kennis,
vaardigheid of houding. Leerresultaten zijn ontwikkeld om een concrete
en consistente interpretatie van het raamwerk mogelijk te maken. \\
\end{longtable}

\emph{Bron: Eigen bewerking door JRC.}

\subsection{Definitie van digitale
competentie}\label{sec-definitie-van-digitale-competentie}

De Raad van de Europese Unie -- Aanbeveling over sleutelcompetenties
voor levenslang leren (Europese Commissie, 2018) definieert digitale
competentie, die toegepast kan worden op zowel kinderen als volwassenen,
als: \ldots het zelfverzekerde, kritische en verantwoorde gebruik van,
en betrokkenheid bij, digitale technologie voor leren, op het werk, en
voor deelname aan de samenleving. Het omvat informatie- en
Data-vaardigheden, communicatie en samenwerking, mediawijsheid, creatie
van digitale content (waaronder programmeren), veiligheid (waaronder
digitaal welzijn en competenties met betrekking tot Cybersecurity),
vraagstukken met betrekking tot intellectueel eigendom,
probleemoplossing en kritisch denken. (European Commission, 2018)

Competenties worden gedefinieerd als een combinatie van kennis,
vaardigheden en houdingen, waarbij:

\begin{itemize}
\item
  \textbf{Kennis} bestaat uit de feiten en cijfers, concepten, ideeën en
  theorieën die reeds vaststaan en het begrip van een bepaald gebied of
  onderwerp ondersteunen;
\item
  \textbf{Vaardigheden} worden gedefinieerd als de bekwaamheid en
  capaciteit om processen uit te voeren en de bestaande kennis te
  gebruiken om resultaten te bereiken; en
\item
  \textbf{Houdingen (attitudes)} beschrijven de instelling en denkwijzen
  om te handelen of te reageren op ideeën, personen of situaties.
\end{itemize}

DigComp 3.0 blijft in lijn met deze algemene definitie, terwijl het
tegelijkertijd recente en opkomende trends en prioriteiten weerspiegelt.

\subsection{Competentiegebieden en
competenties}\label{sec-competentiegebieden-en-competenties}

Figuur~\ref{fig-competentiegebieden-en-competenties} laat zien hoe de
competenties zijn gegroepeerd in de competentiegebieden, terwijl
Tabel~\ref{tbl-comp-geb-beschrijvingen} bovendien de descriptoren voor
elk competentiegebied en elke competentie weergeeft.

\begin{figure}

\pandocbounded{\includegraphics[keepaspectratio]{../output/Vijfhoek_nl.png}}

\caption{\label{fig-competentiegebieden-en-competenties}Competentiegebieden
en competenties van DigComp 3.0}

\end{figure}%

\begin{longtable}[]{@{}
  >{\centering\arraybackslash}p{(\linewidth - 4\tabcolsep) * \real{0.2759}}
  >{\centering\arraybackslash}p{(\linewidth - 4\tabcolsep) * \real{0.2414}}
  >{\centering\arraybackslash}p{(\linewidth - 4\tabcolsep) * \real{0.4828}}@{}}
\caption{Overzicht van gebieden en
competenties}\label{tbl-comp-geb-beschrijvingen}\tabularnewline
\toprule\noalign{}
\begin{minipage}[b]{\linewidth}\centering
GEBIED
\end{minipage} & \begin{minipage}[b]{\linewidth}\centering
TITEL
\end{minipage} & \begin{minipage}[b]{\linewidth}\centering
BESCHRIJVING
\end{minipage} \\
\midrule\noalign{}
\endfirsthead
\toprule\noalign{}
\begin{minipage}[b]{\linewidth}\centering
GEBIED
\end{minipage} & \begin{minipage}[b]{\linewidth}\centering
TITEL
\end{minipage} & \begin{minipage}[b]{\linewidth}\centering
BESCHRIJVING
\end{minipage} \\
\midrule\noalign{}
\endhead
\bottomrule\noalign{}
\endlastfoot
1. INFORMATIE ZOEKEN, EVALUEREN EN BEHEREN & 1.1 Browsen, zoeken en
filteren van informatie & Het formuleren van informatiebehoeften, weten
hoe en waar te zoeken naar informatie en inhoud in digitale omgevingen,
en hiertoe toegang verkrijgen en ertussen navigeren. Geschikte digitale
middelen selecteren om zoekopdrachten in digitale omgevingen te creëren,
uit te voeren en bij te werken, en in staat zijn om onderscheid te maken
tussen relevante en irrelevante informatie en inhoud. \\
& 1.2 Informatie evalueren & De geloofwaardigheid en betrouwbaarheid van
informatiebronnen en inhoud in digitale omgevingen beoordelen en
vergelijken. Informatie en inhoud in digitale omgevingen, en de
processen die worden gebruikt om deze te genereren, interpreteren en
kritisch evalueren. \\
& 1.3 Informatie beheren & Het organiseren, opslaan en ophalen van
informatie en gegevens in digitale omgevingen. Het verzamelen, verwerken
en analyseren van informatie en gegevens in gestructureerde digitale
omgevingen. \\
2. COMMUNICATIE EN SAMENWERKING & 2.1 Interacteren via en met digitale
technologieën & Interacteren via en met een verscheidenheid aan digitale
technologieën, en passende digitale communicatie gebruiken voor een
gegeven context. \\
& 2.2 Delen via digitale technologieën & Informatie en inhoud op een
ethische en verantwoorde manier delen met anderen via passende digitale
technologieën. \\
& 2.3 Betrokkenheid bij burgerschap via digitale technologieën &
Deelnemen aan de samenleving door het ethisch en verantwoord gebruik van
digitale platforms en diensten. Mogelijkheden zoeken voor
zelfontplooiing en participatie via passende digitale technologieën. Je
bewust zijn van je rechten en deze opeisen, en keuzes maken in digitale
omgevingen. \\
& 2.4 Samenwerken via digitale technologieën & Digitale technologieën
ethisch en verantwoord gebruiken voor samenwerking, en voor de
gezamenlijke opbouw en creatie van informatie, bronnen en kennis. \\
& 2.5 Digitaal gedrag & Je bewust zijn van gedragsnormen en weten hoe je
je respectvol moet gedragen bij het gebruik van digitale technologieën
en interactie in digitale omgevingen. Communicatie aanpassen aan
specifieke contexten en aandacht hebben voor culturele, generationele en
andere diversiteit in digitale omgevingen. \\
& 2.6 Digitale identiteit beheren & Het beheren van één of meerdere
digitale identiteiten. Acties ondernemen om je digitale reputatie te
beschermen (hoe je wordt waargenomen op basis van online aanwezigheid)
en je digitale voetafdruk beheren (de gegevens die worden geproduceerd
door het gebruik van digitale platforms en diensten). \\
3. CONTENTCREATIE & 3.1 Digitale inhoud ontwikkelen & Digitale
technologieën op een ethische en verantwoorde manier gebruiken om
uiteenlopende inhoud te creëren en te bewerken. Jezelf uitdrukken via
digitale middelen. \\
& 3.2 Integreren en herwerken van digitale inhoud & Nieuwe informatie en
inhoud aanpassen, verfijnen en integreren in bestaande kennis en bronnen
om nieuwe en originele inhoud en kennis te creëren. \\
& 3.3 Auteursrecht en licenties & Begrijpen hoe auteursrecht en
licenties, evenals de bijbehorende juridische en ethische kwesties, van
toepassing zijn op digitale inhoud, en hoe deze correct moeten worden
toegepast. \\
& 3.4 Computational thinking en programmeren & Stappen begrijpen en
uitvoeren om een probleem te analyseren, deelproblemen te herkennen en
een reeks instructies te plannen en te ontwikkelen voor een
computersysteem om een bepaald probleem op te lossen of een specifieke
taak uit te voeren. \\
4. VEILIGHEID, WELZIJN EN VERANTWOORD GEBRUIK & 4.1 Apparaten beschermen
& Veiligheids- en cybersecuritymaatregelen toepassen om digitale
apparaten en inhoud te beschermen. Je bewust zijn van de veranderende
aard van risico's en bedreigingen in digitale omgevingen, en gepaste
aandacht hebben voor de beveiliging van digitale apparaten en hun
inhoud. \\
& 4.2 Persoonsgegevens en privacy beschermen & Je bewust zijn van je
rechten met betrekking tot persoonsgegevens en privacy in digitale
omgevingen en deze uitoefenen. Privacyrisico's evalueren en beheren en
persoonsgegevens en privacy in digitale omgevingen beschermen. Je eigen
persoonsgegevens en die van anderen op een veilige, ethische en
verantwoorde manier gebruiken en delen. \\
& 4.3 Welzijn ondersteunen & Digitale technologieën gebruiken op
manieren die welzijn en inclusie ondersteunen. Risico's en bedreigingen
voor het fysieke, mentale en sociale welzijn van jezelf en anderen
minimaliseren bij het gebruik van digitale technologieën. Het gebruik
van digitale technologieën in evenwicht brengen met offline activiteiten
om welzijn te ondersteunen. Actie ondernemen om jezelf en anderen te
beschermen tegen mogelijke gevaren in digitale omgevingen (bijv.
cyberpesten, schadelijke inhoud), en weten hoe op dergelijke gevaren te
reageren. \\
& 4.4 Milieueffecten van digitale technologieën & Je bewust zijn van de
milieueffecten van digitale technologieën, waaronder de productie van
apparaten, het gebruik, de reparatie, recycling, verwijdering,
infrastructuur voor gegevensopslag, energieverbruik en het gebruik van
middelen en applicaties. Actie ondernemen om deze impact te beperken en
digitale technologieën gebruiken om duurzaamheid te ondersteunen. \\
5. PROBLEEMIDENTIFICATIE EN -OPLOSSING & 5.1 Technische problemen
identificeren en oplossen & Technische problemen identificeren bij het
bedienen van digitale apparaten en in digitale omgevingen, en deze op
verschillende manieren oplossen. \\
& 5.2 Behoeften en digitale technologische antwoorden identificeren &
Behoeften van jezelf en anderen beoordelen en digitale technologieën
evalueren, selecteren, gebruiken en aanpassen om aan deze behoeften te
voldoen. Digitale omgevingen aanpassen en personaliseren aan de context,
doelen en behoeften (bijv. toegankelijkheid) van jezelf en anderen. \\
& 5.3 Creatieve oplossingen identificeren met behulp van digitale
technologieën & Digitale technologieën gebruiken om verbeteringen aan te
brengen in of nieuwe oplossingen te vinden voor processen en producten,
vanuit een mensgerichte benadering. Individueel en collectief deelnemen
aan kritische denkprocessen en het creatief en doelgericht gebruik van
digitale technologieën, om conceptuele problemen en probleemsituaties te
begrijpen en op te lossen. \\
& 5.4 Behoeften aan digitale competenties identificeren en aanpakken &
Herkennen waar de eigen digitale competentie verbeterd of geactualiseerd
moet worden. Behoeften aan digitale competenties aanpakken binnen een
breder proces van een leven lang leren, en bekwaamheid en autonomie
opbouwen. Anderen ondersteunen bij de ontwikkeling van hun digitale
competenties. Op de hoogte blijven van digitale technologische
ontwikkelingen en de persoonlijke, professionele en maatschappelijke
gevolgen daarvan. \\
\end{longtable}

\emph{Bron: Eigen bewerking door JRC.}

\subsection{Beheersingsniveaus}\label{sec-beheersingsniveaus}

\subsubsection{Beschrijvingen van
beheersingsniveaus}\label{beschrijvingen-van-beheersingsniveaus}

In DigComp 3.0 beschrijven de beheersingsniveaus het niveau van
verwerving van digitale competentie van een individu op basis van een
combinatie van cognitieve inspanning, taakcomplexiteit en mate van
autonomie. DigComp 3.0 maakt onderscheid tussen vier beheersingsniveaus
(Basis, Gemiddeld, Geavanceerd en Zeer Geavanceerd).
Tabel~\ref{tbl-beheersingsniveaus} geeft een korte beschrijving van wat
een individu geacht wordt te weten en te kunnen op elk van de vier
niveaus en de doelen die elk van de vier niveaus dient.

\begin{longtable}[]{@{}
  >{\raggedright\arraybackslash}p{(\linewidth - 4\tabcolsep) * \real{0.2000}}
  >{\raggedright\arraybackslash}p{(\linewidth - 4\tabcolsep) * \real{0.4000}}
  >{\raggedright\arraybackslash}p{(\linewidth - 4\tabcolsep) * \real{0.4000}}@{}}
\caption{DigComp 3.0 Beheersingsniveau beschrijvingen en
doeleinden}\label{tbl-beheersingsniveaus}\tabularnewline
\toprule\noalign{}
\begin{minipage}[b]{\linewidth}\raggedright
Beheersingsniveau
\end{minipage} & \begin{minipage}[b]{\linewidth}\raggedright
Korte beschrijving van competentieverwerving
\end{minipage} & \begin{minipage}[b]{\linewidth}\raggedright
Doel
\end{minipage} \\
\midrule\noalign{}
\endfirsthead
\toprule\noalign{}
\begin{minipage}[b]{\linewidth}\raggedright
Beheersingsniveau
\end{minipage} & \begin{minipage}[b]{\linewidth}\raggedright
Korte beschrijving van competentieverwerving
\end{minipage} & \begin{minipage}[b]{\linewidth}\raggedright
Doel
\end{minipage} \\
\midrule\noalign{}
\endhead
\bottomrule\noalign{}
\endlastfoot
\textbf{Basis} & Op basisniveau onthouden en voeren personen eenvoudige
taken uit met begeleiding waar nodig. & Om persoonlijke, leer- en/of
werkdoelen te ondersteunen en deel te nemen aan de samenleving. \\
\textbf{Gemiddeld} & Op gemiddeld niveau identificeren en voeren
personen welgedefinieerde taken uit en lossen ze welgedefinieerde
problemen autonoom op. & Om persoonlijke, leer- en/of werkdoelen te
ondersteunen en autonoom deel te nemen aan de samenleving. \\
\textbf{Geavanceerd} & Op geavanceerd niveau beoordelen en passen
personen autonoom oplossingen toe op een verscheidenheid aan complexe
taken en passen ze zich aan diverse contexten aan om taken adequaat te
evalueren en uit te voeren, waarbij ze anderen begeleiden indien en
wanneer nodig. & Om persoonlijke, leer- en/of werkdoelen te
ondersteunen, effectief deel te nemen aan de samenleving en anderen te
managen of te ondersteunen bij het bereiken van hun doelen. \\
\textbf{Zeer geavanceerd} & Op zeer geavanceerd niveau beoordelen,
evalueren en lossen personen zeer complexe of gespecialiseerde problemen
op om nieuwe oplossingen te creëren of bestaande aan te passen, waarbij
ze anderen leiden en begeleiden indien en wanneer nodig. & Om
persoonlijke, leer- en/of werkdoelen te ondersteunen, anderen te helpen
effectief deel te nemen aan de samenleving, anderen te leiden of te
ondersteunen bij het bereiken van complexe doelen, en/of leiding te
geven aan of bij te dragen aan verbeteringen in of nieuwe oplossingen
voor zeer complexe problemen. \\
\end{longtable}

\emph{Bron: Eigen bewerking door JRC.}

Elk van de vier beheersingsniveaus van DigComp 3.0 omvat

\begin{enumerate}
\def\labelenumi{\arabic{enumi}.}
\item
  een korte beschrijving van competentieverwerving en
\item
  de doelen die elk niveau kan dienen. Het is belangrijk om te erkennen
  dat de behoeften aan digitale competentie variëren per individu en
  veranderen in de loop van de tijd, als gevolg van levensfasen en
  technologische ontwikkelingen. Een taak is een specifieke activiteit
  waarbij digitale technologie worden gebruikt die bijdragen aan een
  doel, in welke context dan ook -- het dagelijks leven, werk of leren.
  Taken kunnen variëren in omvang, duur en complexiteit en kunnen
  individueel of in samenwerking met anderen worden uitgevoerd. Taken
  kunnen variëren van eenvoudig tot complex. In DigComp 3.0 is een
  eenvoudige taak een taak die goed gedefinieerd is, uit weinig
  onderdelen bestaat en gemakkelijk te begrijpen en te voltooien is. Een
  complexe taak daarentegen is een taak die niet goed gedefinieerd is,
  uit veel verschillende en onderling verbonden onderdelen bestaat en
  daarom ingewikkeld en niet gemakkelijk te begrijpen of te voltooien
  is. De complexiteit van een taak wordt vaak alleen beschreven op basis
  van taakkenmerken, maar de ervaring en kenmerken van een individu dat
  een taak uitvoert, zijn ook relevant (Chen et al., 2023). Op de
  niveaus Geavanceerd en Zeer Geavanceerd van digitale competentie
  kunnen individuen naast digitale competenties ook putten uit ervaring
  en/of gespecialiseerde kennis om taken te voltooien. Hier verwijzen
  gespecialiseerde kennis en vaardigheden naar competenties die
  betrekking hebben op een bepaalde discipline, vakgebied of terrein.
\end{enumerate}

\newpage{}

\subsection{Leerresultaten}\label{sec-leerresultaten}

\subsubsection{1.1 Browsen, zoeken en filteren van
informatie}\label{browsen-zoeken-en-filteren-van-informatie}

\begin{longtable}[]{@{}
  >{\centering\arraybackslash}p{(\linewidth - 8\tabcolsep) * \real{0.1053}}
  >{\centering\arraybackslash}p{(\linewidth - 8\tabcolsep) * \real{0.3947}}
  >{\centering\arraybackslash}p{(\linewidth - 8\tabcolsep) * \real{0.2105}}
  >{\centering\arraybackslash}p{(\linewidth - 8\tabcolsep) * \real{0.1842}}
  >{\centering\arraybackslash}p{(\linewidth - 8\tabcolsep) * \real{0.1053}}@{}}
\toprule\noalign{}
\begin{minipage}[b]{\linewidth}\centering
ID
\end{minipage} & \begin{minipage}[b]{\linewidth}\centering
Leerresultaat
\end{minipage} & \begin{minipage}[b]{\linewidth}\centering
Niveau
\end{minipage} & \begin{minipage}[b]{\linewidth}\centering
K/V/H
\end{minipage} & \begin{minipage}[b]{\linewidth}\centering
AI
\end{minipage} \\
\midrule\noalign{}
\endhead
\bottomrule\noalign{}
\endlastfoot
LO1.1.01 & De voordelen inzien van het gebruik van verschillende
digitale zoekhulpmiddelen en -methoden, afhankelijk van het doel. &
Basis & H & AI-Impliciet \\
LO1.1.02 & Herkennen dat de resultaten of outputs van een zoekopdracht
afhankelijk zijn van het gebruikte digitale zoekhulpmiddel en de manier
waarop een individu de zoekopdracht specificeert. & Basis & K &
AI-Impliciet \\
LO1.1.03 & Herkennen dat zoekresultaten of outputs informatie kunnen
bevatten die mogelijk niet relevant is. & Basis & K & AI-Impliciet \\
LO1.1.04 & De belangrijkste kenmerken identificeren van veelgebruikte
AI-gestuurde en traditionele digitale zoekhulpmiddelen. & Basis & K &
AI-Expliciet \\
LO1.1.05 & Digitale zoekhulpmiddelen gebruiken om eenvoudige
informatiezoekopdrachten uit te voeren. & Basis & V & AI-Impliciet \\
LO1.1.06 & Digitale zoekhulpmiddelen gebruiken om bestaande resultaten
of outputs te verfijnen of bij te werken. & Basis & V & AI-Impliciet \\
LO1.1.07 & Doelgericht nieuwe digitale zoekhulpmiddelen en
zoekfunctionaliteiten verkennen. & Gemiddeld & H & AI-Impliciet \\
LO1.1.08 & Strategieën identificeren die relevantere digitale
zoekresultaten of outputs opleveren. & Gemiddeld & K & AI-Impliciet \\
LO1.1.09 & Onderscheid maken tussen meer en minder relevante digitale
zoekresultaten of outputs. & Gemiddeld & K & AI-Impliciet \\
LO1.1.10 & Passende digitale zoekhulpmiddelen selecteren op basis van
informatiebehoeften. & Gemiddeld & V & AI-Impliciet \\
LO1.1.11 & Informatiebehoeften vertalen naar effectieve digitale
zoekopdrachten, commando's of verklaringen. & Gemiddeld & V &
AI-Impliciet \\
LO1.1.12 & Passende strategieën toepassen om bestaande digitale
resultaten of outputs te verfijnen of te filteren. & Gemiddeld & V &
AI-Impliciet \\
LO1.1.13 & Voortdurend functies en kenmerken van bekende en nieuwe
digitale zoekhulpmiddelen verkennen. & Geavanceerd & H & AI-Impliciet \\
LO1.1.14 & Prioriteit geven aan het verdiepen van de eigen bestaande
zoekvaardigheden. & Geavanceerd & H & AI-Impliciet \\
LO1.1.15 & Een verscheidenheid aan digitale zoekhulpmiddelen en
strategieën combineren om aan complexe informatiebehoeften te voldoen. &
Geavanceerd & V & AI-Impliciet \\
LO1.1.16 & Anderen helpen bij het ontwikkelen van hun digitale
zoekvaardigheden. & Geavanceerd & V & AI-Impliciet \\
LO1.1.17 & Op de hoogte blijven van ontwikkelingen in digitale
zoektechnologieën. & Zeer geavanceerd & H & AI-Impliciet \\
LO1.1.18 & Ontwikkelingen in digitale zoektechnologieën in een gegeven
context beoordelen en evalueren ter ondersteuning van besluitvorming. &
Zeer geavanceerd & V & AI-Impliciet \\
LO1.1.19 & Een verscheidenheid aan digitale zoekhulpmiddelen en
strategieën combineren om aan zeer complexe of gespecialiseerde
informatiebehoeften te voldoen. & Zeer geavanceerd & V & AI-Impliciet \\
LO1.1.20 & Anderen helpen bij het uitvoeren en verfijnen van complexe of
gespecialiseerde zoekopdrachten in digitale omgevingen. & Zeer
geavanceerd & V & AI-Impliciet \\
LO1.1.21 & Bijdragen aan verbeteringen in of nieuwe oplossingen voor
complexe of gespecialiseerde zoekopdrachten in digitale omgevingen. &
Zeer geavanceerd & V & AI-Impliciet \\
\end{longtable}

\subsubsection{1.2 Informatie evalueren}\label{informatie-evalueren}

\begin{longtable}[]{@{}
  >{\centering\arraybackslash}p{(\linewidth - 8\tabcolsep) * \real{0.1053}}
  >{\centering\arraybackslash}p{(\linewidth - 8\tabcolsep) * \real{0.3947}}
  >{\centering\arraybackslash}p{(\linewidth - 8\tabcolsep) * \real{0.2105}}
  >{\centering\arraybackslash}p{(\linewidth - 8\tabcolsep) * \real{0.1842}}
  >{\centering\arraybackslash}p{(\linewidth - 8\tabcolsep) * \real{0.1053}}@{}}
\toprule\noalign{}
\begin{minipage}[b]{\linewidth}\centering
ID
\end{minipage} & \begin{minipage}[b]{\linewidth}\centering
Leerresultaat
\end{minipage} & \begin{minipage}[b]{\linewidth}\centering
Niveau
\end{minipage} & \begin{minipage}[b]{\linewidth}\centering
K/V/H
\end{minipage} & \begin{minipage}[b]{\linewidth}\centering
AI
\end{minipage} \\
\midrule\noalign{}
\endhead
\bottomrule\noalign{}
\endlastfoot
LO1.2.01 & De voordelen inzien van een voorzichtige aanpak bij het
interpreteren van informatie en inhoud in digitale omgevingen. & Basis &
H & - \\
LO1.2.02 & Onderscheid maken tussen de bron van digitale inhoud en de
digitale inhoud zelf. & Basis & K & AI-Impliciet \\
LO1.2.03 & Herkennen dat sommige digitale informatiebronnen en systemen
mogelijk niet betrouwbaar zijn. & Basis & K & AI-Impliciet \\
LO1.2.04 & Herkennen dat het moeilijk kan zijn om onderscheid te maken
tussen informatie en inhoud gegenereerd door mensen en door AI-systemen.
& Basis & K & AI-Expliciet \\
LO1.2.05 & Voorbeelden herkennen van misinformatie, desinformatie en
bronnen van vooringenomenheid (bias). & Basis & K & AI-Impliciet \\
LO1.2.06 & Voorbeelden herkennen van beïnvloeding via sociale media en
filterbubbels. & Basis & K & AI-Impliciet \\
LO1.2.07 & Een basisbeoordeling maken van de betrouwbaarheid en
geloofwaardigheid van digitale informatiebronnen en inhoud. & Basis & V
& AI-Impliciet \\
LO1.2.08 & Het nut onderkennen van het kritisch beoordelen van de
geloofwaardigheid en betrouwbaarheid van informatie, inhoud en hun
bronnen in digitale omgevingen. & Gemiddeld & H & AI-Impliciet \\
LO1.2.09 & Mogelijke gevolgen herkennen van misinformatie en
desinformatie in digitale omgevingen voor zichzelf en anderen. &
Gemiddeld & K & AI-Impliciet \\
LO1.2.10 & Methoden beschrijven om de bron van online gevonden
informatie te identificeren. & Gemiddeld & K & AI-Impliciet \\
LO1.2.11 & De doelen van factcheck-diensten definiëren. & Gemiddeld & K
& AI-Impliciet \\
LO1.2.12 & De concepten en doelen van pre-bunking en de-bunking
herkennen in digitale contexten. & Gemiddeld & K & AI-Impliciet \\
LO1.2.13 & Herkenen dat de data die ingezet wordt voor het trainen van
AI-systemen en de wijze waarop zij getraind worden van invloed zijn op
de betrouwbaarheid van de informatie die zij leveren. & Gemiddeld & K &
AI-Expliciet \\
LO1.2.14 & Herkennen dat sommige digitale technologieën, zoals
AI-systemen, kunnen functioneren als een `black box', waardoor het
moeilijk is om uit te leggen waarom of hoe een output is geproduceerd. &
Gemiddeld & K & AI-Expliciet \\
LO1.2.15 & Voorbeelden identificeren van menselijke (cognitieve,
affectieve) bias en AI-systeem (data, training) bias in relatie tot het
genereren en interpreteren van informatie. & Gemiddeld & K &
AI-Expliciet \\
LO1.2.16 & Herkennen dat AI-systemen output kunnen produceren die
onnauwkeurig is, zelfs als deze aannemelijk lijkt. & Gemiddeld & K &
AI-Expliciet \\
LO1.2.17 & Herkennen dat de mensen die een AI-systeem gebruiken
verantwoordelijk zijn voor het controleren van de kwaliteit en
geldigheid van de gegenereerde informatie en inhoud. & Gemiddeld & K &
AI-Expliciet \\
LO1.2.18 & De aanwezigheid herkennen van gebruikerssturende strategieën
in digitale omgevingen, zoals clickbait, nudging en gamificatie. &
Gemiddeld & K & AI-Impliciet \\
LO1.2.19 & Pre-bunking en de-bunking strategieën toepassen om
onbetrouwbare bronnen en inhoud in digitale omgevingen te verwerpen of
in diskrediet te brengen. & Gemiddeld & V & AI-Impliciet \\
LO1.2.20 & Effectief reageren op gebruikerssturende strategieën in
digitale omgevingen, zoals clickbait, nudging en gamificatie. &
Gemiddeld & V & AI-Impliciet \\
LO1.2.21 & Kritisch de betrouwbaarheid beoordelen van bronnen,
informatie en inhoud in digitale omgevingen, rekening houdend met de rol
van AI-systemen, personalisatie-effecten en commerciële of andere
belangen. & Gemiddeld & V & AI-Expliciet \\
LO1.2.22 & Voortdurend onderzoeken hoe AI-systemen, vooroordelen en
verschillende belangen de generatie, presentatie en interpretatie van
informatie en inhoud in digitale omgevingen vormgeven. & Geavanceerd & H
& AI-Expliciet \\
LO1.2.23 & Persoonlijke, sociale en politieke gevolgen beschrijven van
misinformatie, desinformatie, bronnen van bias, beïnvloeding via sociale
media en filterbubbels. & Geavanceerd & K & AI-Impliciet \\
LO1.2.24 & Kenmerken van betrouwbare digitale technologieën, zoals
AI-systemen, beschrijven. & Geavanceerd & K & AI-Expliciet \\
LO1.2.25 & Methoden beschrijven die gebruikt kunnen worden om deepfakes
te identificeren. & Geavanceerd & K & AI-Expliciet \\
LO1.2.26 & Grondig de betrouwbaarheid en nauwkeurigheid beoordelen van
een diversiteit aan bronnen, informatie en inhoud in digitale
omgevingen, rekening houdend met een reeks mogelijke beïnvloedende
factoren. & Geavanceerd & V & AI-Impliciet \\
LO1.2.27 & Anderen ondersteunen bij het ontwikkelen van hun vaardigheden
om de betrouwbaarheid van bronnen, informatie en inhoud in digitale
omgevingen te beoordelen. & Geavanceerd & V & AI-Impliciet \\
LO1.2.28 & Het ontwikkelen van veerkracht tegen misinformatie en
desinformatie in digitale omgevingen onder individuen en/of groepen
bevorderen en ondersteunen. & Zeer geavanceerd & H & AI-Impliciet \\
LO1.2.29 & Systematisch bronnen, informatie en inhoud in digitale
omgevingen beoordelen en evalueren ter ondersteuning van complexe
besluitvorming. & Zeer geavanceerd & V & AI-Impliciet \\
LO1.2.30 & Anderen helpen bij het ontwikkelen van vaardigheden om
informatie en inhoud in digitale omgevingen kritisch te evalueren. &
Zeer geavanceerd & V & AI-Impliciet \\
LO1.2.31 & Leiding geven aan of bijdragen aan initiatieven die een
nauwkeurige interpretatie van data, informatie en inhoud in digitale
omgevingen ondersteunen. & Zeer geavanceerd & V & AI-Impliciet \\
\end{longtable}

\subsubsection{1.3 Informatie beheren}\label{informatie-beheren}

\begin{longtable}[]{@{}
  >{\centering\arraybackslash}p{(\linewidth - 8\tabcolsep) * \real{0.1053}}
  >{\centering\arraybackslash}p{(\linewidth - 8\tabcolsep) * \real{0.3947}}
  >{\centering\arraybackslash}p{(\linewidth - 8\tabcolsep) * \real{0.2105}}
  >{\centering\arraybackslash}p{(\linewidth - 8\tabcolsep) * \real{0.1842}}
  >{\centering\arraybackslash}p{(\linewidth - 8\tabcolsep) * \real{0.1053}}@{}}
\toprule\noalign{}
\begin{minipage}[b]{\linewidth}\centering
ID
\end{minipage} & \begin{minipage}[b]{\linewidth}\centering
Leerresultaat
\end{minipage} & \begin{minipage}[b]{\linewidth}\centering
Niveau
\end{minipage} & \begin{minipage}[b]{\linewidth}\centering
K/V/H
\end{minipage} & \begin{minipage}[b]{\linewidth}\centering
AI
\end{minipage} \\
\midrule\noalign{}
\endhead
\bottomrule\noalign{}
\endlastfoot
LO1.3.01 & De voordelen inzien van het beheren en organiseren van
informatie in digitale omgevingen. & Basis & H & - \\
LO1.3.02 & Functies van dataverwijdering, -herstel en -back-up
herkennen. & Basis & K & - \\
LO1.3.03 & De belangrijkste eigenschappen van digitale bestanden en
mappen identificeren. & Basis & K & - \\
LO1.3.04 & In algemene termen het concept van data herkennen. & Basis &
K & - \\
LO1.3.05 & Digitale bestanden downloaden, opslaan, ophalen, verplaatsen
en verwijderen. & Basis & V & - \\
LO1.3.06 & Eenvoudige data organiseren en formatteren in een
gestructureerde digitale omgeving, zoals in spreadsheets. & Basis & V &
- \\
LO1.3.07 & Eigen contacten bijwerken, bijvoorbeeld op de telefoon, in
e-mail of op sociale media. & Basis & V & - \\
LO1.3.08 & Het belang inzien van zorgvuldig en ethisch beheer van data
en informatie in digitale omgevingen. & Gemiddeld & H & AI-Impliciet \\
LO1.3.09 & Herkennen dat digitale bestanden en mappen kunnen worden
(her)noemd en georganiseerd op een door de gebruiker gewenste manier. &
Gemiddeld & K & - \\
LO1.3.10 & Herkennen dat digitale bestanden op verschillende locaties
kunnen worden opgeslagen (apparaten, externe opslag en clouddiensten) en
van de ene naar de andere locatie kunnen worden overgezet. & Gemiddeld &
K & - \\
LO1.3.11 & Veelvoorkomende hulpmiddelen voor gegevensverzameling en hun
belangrijkste functies identificeren. & Gemiddeld & K & AI-Impliciet \\
LO1.3.12 & Verantwoordelijkheden definiëren die verbonden zijn aan het
gebruik van hulpmiddelen voor gegevensverzameling. & Gemiddeld & K &
- \\
LO1.3.13 & Veelvoorkomende soorten data en hun formaten identificeren. &
Gemiddeld & K & - \\
LO1.3.14 & Naamgevingsconventies toepassen op digitale bestanden en
hiërarchieën op digitale mappen. & Gemiddeld & V & - \\
LO1.3.15 & Bestanden beheren, opslaan en verwijderen op digitale
apparaten, externe opslag en clouddiensten. & Gemiddeld & V & - \\
LO1.3.16 & Informatie beheren in de eigen digitale accounts, zoals
e-mail. & Gemiddeld & V & - \\
LO1.3.17 & Hulpmiddelen voor gegevensverzameling gebruiken voor
eenvoudige verwerking van data en informatie, zoals quizzen, peilingen
of enquêtes. & Gemiddeld & V & AI-Impliciet \\
LO1.3.18 & Data organiseren en formatteren in een gestructureerde
digitale omgeving, zoals in spreadsheets. & Gemiddeld & V & - \\
LO1.3.19 & Eenvoudige formules toepassen op data in een gestructureerde
digitale omgeving, zoals in spreadsheets. & Gemiddeld & V & - \\
LO1.3.20 & Prioriteit geven aan ethisch en transparant beheer en
verwerking van data en informatie in digitale omgevingen. & Geavanceerd
& H & AI-Impliciet \\
LO1.3.21 & Rekening houden met mogelijke bronnen van fouten of
onnauwkeurigheden bij het beheer en de verwerking van data en informatie
in digitale omgevingen. & Geavanceerd & H & AI-Impliciet \\
LO1.3.22 & Mogelijke bronnen van fouten of onnauwkeurigheden in
informatie of data in digitale omgevingen identificeren. & Geavanceerd &
K & AI-Impliciet \\
LO1.3.23 & De belangrijkste stappen beschrijven bij het beheren,
verwerken en analyseren van informatie en data in digitale omgevingen. &
Geavanceerd & K & AI-Impliciet \\
LO1.3.24 & Kenmerken van open data beschrijven (voorbeelden,
toepassingen, voordelen en beperkingen). & Geavanceerd & K &
AI-Impliciet \\
LO1.3.25 & Kenmerken van big data beschrijven (voorbeelden,
toepassingen, voordelen en beperkingen). & Geavanceerd & K &
AI-Impliciet \\
LO1.3.26 & Verschillende functies toepassen om data en informatie in
digitale omgevingen over te dragen en te beheren. & Geavanceerd & V &
AI-Impliciet \\
LO1.3.27 & Een reeks digitale hulpmiddelen en methoden gebruiken om een
verscheidenheid aan data en informatie te verzamelen en te verwerken. &
Geavanceerd & V & AI-Impliciet \\
LO1.3.28 & Passende analyses van informatie en data in digitale
omgevingen toepassen om bij te dragen aan complexe besluitvorming. &
Geavanceerd & V & AI-Impliciet \\
LO1.3.29 & Anderen helpen met het beheer, de verwerking en de analyse
van data en informatie in digitale omgevingen. & Geavanceerd & V &
AI-Impliciet \\
LO1.3.30 & Het belang inzien van het structureren en documenteren van
data en informatie in digitale omgevingen ten behoeve van anderen. &
Zeer geavanceerd & H & - \\
LO1.3.31 & Op de hoogte blijven van technologische ontwikkelingen op het
gebied van data- en informatiebeheer en -analyse. & Zeer geavanceerd & H
& AI-Impliciet \\
LO1.3.32 & Strategieën ontwikkelen en implementeren voor complex of
gespecialiseerd beheer, verwerking en analyse van data en informatie in
digitale omgevingen. & Zeer geavanceerd & V & AI-Impliciet \\
LO1.3.33 & Verschillende hulpmiddelen en methoden zoals big
data-technieken of simulaties gebruiken om complexe data of grote
hoeveelheden informatie te verwerken, te beheren of te analyseren. &
Zeer geavanceerd & V & AI-Impliciet \\
LO1.3.34 & Leiding geven aan of bijdragen aan initiatieven die anderen
ondersteunen bij geavanceerd informatie- en databeheer, -verwerking en
-analyse in digitale omgevingen. & Zeer geavanceerd & V &
AI-Impliciet \\
LO1.3.35 & Bijdragen aan verbeteringen in of nieuwe oplossingen voor
complex databeheer, -verwerking of -analyse in digitale omgevingen. &
Zeer geavanceerd & V & AI-Impliciet \\
\end{longtable}

\subsubsection{2.1 Interacteren via en met digitale
technologieën}\label{interacteren-via-en-met-digitale-technologieuxebn}

\begin{longtable}[]{@{}
  >{\centering\arraybackslash}p{(\linewidth - 8\tabcolsep) * \real{0.1053}}
  >{\centering\arraybackslash}p{(\linewidth - 8\tabcolsep) * \real{0.3947}}
  >{\centering\arraybackslash}p{(\linewidth - 8\tabcolsep) * \real{0.2105}}
  >{\centering\arraybackslash}p{(\linewidth - 8\tabcolsep) * \real{0.1842}}
  >{\centering\arraybackslash}p{(\linewidth - 8\tabcolsep) * \real{0.1053}}@{}}
\toprule\noalign{}
\begin{minipage}[b]{\linewidth}\centering
ID
\end{minipage} & \begin{minipage}[b]{\linewidth}\centering
Leerresultaat
\end{minipage} & \begin{minipage}[b]{\linewidth}\centering
Niveau
\end{minipage} & \begin{minipage}[b]{\linewidth}\centering
K/V/H
\end{minipage} & \begin{minipage}[b]{\linewidth}\centering
AI
\end{minipage} \\
\midrule\noalign{}
\endhead
\bottomrule\noalign{}
\endlastfoot
LO2.1.01 & Het belang inzien van het rekening houden met de voorkeuren
van anderen in digitale communicatie. & Basis & H & - \\
LO2.1.02 & Onderscheid maken tussen synchrone en asynchrone vormen van
digitale communicatie. & Basis & K & - \\
LO2.1.03 & Verschillen identificeren tussen digitale en niet-digitale
interacties. & Basis & K & AI-Impliciet \\
LO2.1.04 & Onderscheid maken tussen fysieke en virtuele realiteiten. &
Basis & K & AI-Impliciet \\
LO2.1.05 & Basiskenmerken en -functies van digitale
communicatiehulpmiddelen identificeren. & Basis & K & AI-Impliciet \\
LO2.1.06 & Basiskenmerken identificeren van virtuele assistenten
(chatbots) en AI-systemen die worden gebruikt in communicatiecontexten.
& Basis & K & AI-Expliciet \\
LO2.1.07 & Belangrijke verschillen herkennen tussen mens-machine
interacties en mens-tot-mens interacties. & Basis & K & AI-Impliciet \\
LO2.1.08 & In algemene termen herkennen wat een robot is, inclusief de
niet-menselijke aard ervan. & Basis & K & - \\
LO2.1.09 & Herkennen dat mensen interageren met robots om taken uit te
voeren. & Basis & K & - \\
LO2.1.10 & Basiskenmerken van digitale communicatiehulpmiddelen
gebruiken om te interageren met individuen en groepen. & Basis & V &
AI-Impliciet \\
LO2.1.11 & Het belang inzien van het afstemmen van de eigen digitale
communicatie op specifieke contexten. & Gemiddeld & H & - \\
LO2.1.12 & Herkennen dat er een realiteit-virtualiteit-continuüm bestaat
in digitale omgevingen. & Gemiddeld & K & AI-Impliciet \\
LO2.1.13 & De belangrijkste kenmerken en functies beschrijven van een
reeks digitale communicatiehulpmiddelen. & Gemiddeld & K &
AI-Impliciet \\
LO2.1.14 & Voordelen en beperkingen beschrijven van virtuele assistenten
(chatbots) en AI-systemen in digitale communicatiecontexten. & Gemiddeld
& K & AI-Expliciet \\
LO2.1.15 & Contexten identificeren waarin asynchrone of synchrone
digitale communicatie, of niet-digitale communicatie, het beste werkt. &
Gemiddeld & K & - \\
LO2.1.16 & Belangrijke kenmerken van robots identificeren (zoals
sensoren, software, bewegingsregelaars en menselijke interface). &
Gemiddeld & K & AI-Impliciet \\
LO2.1.17 & Voorbeelden definiëren van hoe mensen met robots kunnen
interageren. & Gemiddeld & K & AI-Impliciet \\
LO2.1.18 & Herkennen dat robots kunnen werken met verschillende mate van
autonomie. & Gemiddeld & K & AI-Impliciet \\
LO2.1.19 & Geschikte communicatiemiddelen en -hulpmiddelen selecteren,
rekening houdend met digitale en niet-digitale opties, voor een gegeven
context of doel. & Gemiddeld & V & AI-Impliciet \\
LO2.1.20 & Vragen, commando's of verklaringen (prompts) ontwikkelen en
verfijnen voor virtuele assistenten (chatbots) en AI-systemen ter
ondersteuning van niet-complexe digitale interacties. & Gemiddeld & V &
AI-Expliciet \\
LO2.1.21 & Meerdere kenmerken van verschillende digitale
communicatiehulpmiddelen gebruiken om te interageren met en het beheren
van individuen, groepen en kanalen. & Gemiddeld & V & AI-Impliciet \\
LO2.1.22 & Communicatie in digitale omgevingen voortdurend aanpassen als
reactie op een verscheidenheid aan contexten. & Geavanceerd & H & - \\
LO2.1.23 & Digitale communicatiehulpmiddelen en -methoden combineren
voor complexe communicatietaken. & Geavanceerd & V & AI-Impliciet \\
LO2.1.24 & Systematisch vragen, commando's of verklaringen (prompts)
voor AI-systemen ontwikkelen en geleidelijk verfijnen om complexe
digitale interacties af te handelen. & Geavanceerd & V & AI-Expliciet \\
LO2.1.25 & Voordelen en nadelen van robottoepassingen in een specifieke
context beoordelen. & Geavanceerd & V & AI-Impliciet \\
LO2.1.26 & Anderen helpen bij het beoordelen en selecteren van geschikte
digitale communicatiehulpmiddelen voor een bepaald doel. & Geavanceerd &
V & AI-Impliciet \\
LO2.1.27 & Organiseren en/of modereren van complexe digitale
evenementen. & Geavanceerd & V & AI-Impliciet \\
LO2.1.28 & Op de hoogte blijven van ontwikkelingen in hulpmiddelen en
methoden voor digitale communicatie en interactie. & Zeer geavanceerd &
H & AI-Impliciet \\
LO2.1.29 & Digitale communicatie- en interactiehulpmiddelen beoordelen
en combineren voor zeer complexe of nieuwe taken. & Zeer geavanceerd & V
& AI-Impliciet \\
LO2.1.30 & Begeleiding, ondersteuning of leiderschap bieden bij het
geavanceerde gebruik van digitale communicatie- en
interactiehulpmiddelen. & Zeer geavanceerd & V & AI-Impliciet \\
LO2.1.31 & Leiden van of bijdragen aan verbeteringen in of nieuwe
oplossingen voor digitale communicatie of mens-machine interactie. &
Zeer geavanceerd & V & AI-Impliciet \\
\end{longtable}

\subsubsection{2.2 Delen via digitale
technologieën}\label{delen-via-digitale-technologieuxebn}

\begin{longtable}[]{@{}
  >{\centering\arraybackslash}p{(\linewidth - 8\tabcolsep) * \real{0.1053}}
  >{\centering\arraybackslash}p{(\linewidth - 8\tabcolsep) * \real{0.3947}}
  >{\centering\arraybackslash}p{(\linewidth - 8\tabcolsep) * \real{0.2105}}
  >{\centering\arraybackslash}p{(\linewidth - 8\tabcolsep) * \real{0.1842}}
  >{\centering\arraybackslash}p{(\linewidth - 8\tabcolsep) * \real{0.1053}}@{}}
\toprule\noalign{}
\begin{minipage}[b]{\linewidth}\centering
ID
\end{minipage} & \begin{minipage}[b]{\linewidth}\centering
Leerresultaat
\end{minipage} & \begin{minipage}[b]{\linewidth}\centering
Niveau
\end{minipage} & \begin{minipage}[b]{\linewidth}\centering
K/V/H
\end{minipage} & \begin{minipage}[b]{\linewidth}\centering
AI
\end{minipage} \\
\midrule\noalign{}
\endhead
\bottomrule\noalign{}
\endlastfoot
LO2.2.01 & Het belang inzien van ethisch en verantwoord delen van
informatie en inhoud in digitale omgevingen. & Basis & H &
AI-Impliciet \\
LO2.2.02 & Voordelen en risico's herkennen van het delen van informatie
en inhoud in de digitale omgeving. & Basis & K & AI-Impliciet \\
LO2.2.03 & Functies en toepassingen van sociale media identificeren,
evenals voorbeelden van veelvoorkomende sociale mediaplatforms. & Basis
& K & AI-Impliciet \\
LO2.2.04 & Herkennen dat individuen kunnen kiezen hoe en wat ze willen
delen via digitale technologieën. & Basis & K & - \\
LO2.2.05 & Herkennen dat informatie en inhoud - waarvan niet alles waar
of nauwkeurig is - op verschillende manieren kan worden gedeeld door
zowel AI-systemen als mensen. & Basis & K & AI-Expliciet \\
LO2.2.06 & Doel en doelgroep identificeren van informatie en inhoud die
gedeeld gaat worden in digitale omgevingen. & Basis & K & - \\
LO2.2.07 & Eenvoudige processen gebruiken om informatie en inhoud op een
passende manier en in overeenstemming met de doelen te delen in digitale
omgevingen. & Basis & V & - \\
LO2.2.08 & Het belang inzien van het beoordelen van de waarde en
nauwkeurigheid van informatie en inhoud voordat deze in digitale
omgevingen wordt gedeeld. & Gemiddeld & H & - \\
LO2.2.09 & Verantwoordelijkheden definiëren die verbonden zijn aan het
delen van informatie en inhoud in digitale omgevingen. & Gemiddeld & K &
AI-Impliciet \\
LO2.2.10 & Effectieve en ethische manieren beschrijven om informatie en
inhoud te delen in verschillende digitale omgevingen. & Gemiddeld & K &
AI-Impliciet \\
LO2.2.11 & Mogelijke risico's, voordelen en ethische overwegingen
beoordelen van het delen van informatie en inhoud in verschillende
digitale omgevingen. & Gemiddeld & V & AI-Impliciet \\
LO2.2.12 & Informatie en inhoud op een effectieve en ethische manier
delen in verschillende digitale omgevingen. & Gemiddeld & V &
AI-Impliciet \\
LO2.2.13 & Misinformatie en desinformatie die gedeeld is in digitale
omgevingen rapporteren of markeren. & Gemiddeld & V & AI-Impliciet \\
LO2.2.14 & De waarde inzien van het delen van digitale informatie en
inhoud om anderen te helpen, bijvoorbeeld via Open Educational Resources
(OER). & Geavanceerd & H & - \\
LO2.2.15 & Informatie en inhoud op een effectieve en ethische manier
delen in digitale omgevingen ter ondersteuning van persoonlijke, leer-
of professionele doelen van zichzelf en anderen. & Geavanceerd & V &
AI-Impliciet \\
LO2.2.16 & Anderen adviseren over effectieve en ethische manieren om
informatie en inhoud te delen in digitale omgevingen. & Geavanceerd & V
& AI-Impliciet \\
LO2.2.17 & Nieuwe en alternatieve middelen verkennen voor het complex
delen van informatie en inhoud in digitale omgevingen. & Zeer
geavanceerd & H & AI-Impliciet \\
LO2.2.18 & Complex delen van informatie en inhoud over verschillende
digitale technologieën heen faciliteren. & Zeer geavanceerd & V &
AI-Impliciet \\
LO2.2.19 & Bijdragen aan complexe of gespecialiseerde initiatieven voor
het delen van informatie en inhoud in digitale omgevingen. & Zeer
geavanceerd & V & AI-Impliciet \\
LO2.2.20 & Leiden van of bijdragen aan verbeteringen in of nieuwe
oplossingen voor het delen van complexe informatie en inhoud in digitale
omgevingen. & Zeer geavanceerd & V & AI-Impliciet \\
\end{longtable}

\subsubsection{2.3 Betrokkenheid bij burgerschap via digitale
technologieën}\label{betrokkenheid-bij-burgerschap-via-digitale-technologieuxebn}

\begin{longtable}[]{@{}
  >{\centering\arraybackslash}p{(\linewidth - 8\tabcolsep) * \real{0.1053}}
  >{\centering\arraybackslash}p{(\linewidth - 8\tabcolsep) * \real{0.3947}}
  >{\centering\arraybackslash}p{(\linewidth - 8\tabcolsep) * \real{0.2105}}
  >{\centering\arraybackslash}p{(\linewidth - 8\tabcolsep) * \real{0.1842}}
  >{\centering\arraybackslash}p{(\linewidth - 8\tabcolsep) * \real{0.1053}}@{}}
\toprule\noalign{}
\begin{minipage}[b]{\linewidth}\centering
ID
\end{minipage} & \begin{minipage}[b]{\linewidth}\centering
Leerresultaat
\end{minipage} & \begin{minipage}[b]{\linewidth}\centering
Niveau
\end{minipage} & \begin{minipage}[b]{\linewidth}\centering
K/V/H
\end{minipage} & \begin{minipage}[b]{\linewidth}\centering
AI
\end{minipage} \\
\midrule\noalign{}
\endhead
\bottomrule\noalign{}
\endlastfoot
LO2.3.01 & De potentiële voordelen inzien van digitale technologieën
voor de eigen empowerment en participatie en die van anderen. & Basis &
H & - \\
LO2.3.02 & Digitaal burgerschap herkennen als het vermogen om actief en
verantwoordelijk deel te nemen aan gemeenschappen door betrokkenheid bij
digitale technologieën. & Basis & K & - \\
LO2.3.03 & Voorbeelden identificeren van maatschappelijke participatie
in digitale omgevingen. & Basis & K & - \\
LO2.3.04 & Herkennen dat digitale technologieën - inclusief
moeilijkheden bij de toegang ertoe - bepaalde groepen of individuen
kunnen uitsluiten. & Basis & K & AI-Impliciet \\
LO2.3.05 & Herkennen dat er wetten en regels zijn om de rechten van
gebruikers van digitale platforms en diensten te beschermen. & Basis & K
& AI-Impliciet \\
LO2.3.06 & De belangrijkste doelen en functies identificeren van
(persoonlijk relevante) digitale platforms en diensten. & Basis & K &
AI-Impliciet \\
LO2.3.07 & Herkennen dat individuen een actieve rol kunnen spelen bij
het beoordelen of verbeteren van online producten en diensten. & Basis &
K & - \\
LO2.3.08 & Digitale hulpmiddelen gebruiken om te zoeken naar en het
vinden van gemeenschappen voor maatschappelijke participatie over
onderwerpen van belang. & Basis & V & AI-Impliciet \\
LO2.3.09 & (Persoonlijk relevante) digitale platforms en diensten
gebruiken, indien nodig hulp zoekend. & Basis & V & AI-Impliciet \\
LO2.3.10 & Prioriteit geven aan het verkennen van manieren waarop
digitale technologieën de eigen maatschappelijke participatie kunnen
versterken. & Gemiddeld & H & AI-Impliciet \\
LO2.3.11 & Het belang inzien van het identificeren van uitgesloten of
gemarginaliseerde mensen en groepen in digitale omgevingen. & Gemiddeld
& H & AI-Impliciet \\
LO2.3.12 & Deelnemen aan discussies over onderwerpen die te maken hebben
met digitaal burgerschap. & Gemiddeld & H & AI-Impliciet \\
LO2.3.13 & Herkennen dat digitale participatie de actieve betrokkenheid
bij de samenleving is door het gebruik van digitale technologieën. &
Gemiddeld & K & - \\
LO2.3.14 & Herkennen dat maatschappelijke participatie plaatsvindt op
een continuüm. & Gemiddeld & K & - \\
LO2.3.15 & Belangrijke vrijheden, rechten en verantwoordelijkheden van
individuen herkennen onder relevante digitale wetten en regels. &
Gemiddeld & K & AI-Impliciet \\
LO2.3.16 & Definiëren hoe belangrijke rechten uit te oefenen in digitale
omgevingen. & Gemiddeld & K & AI-Impliciet \\
LO2.3.17 & Het concept van de platformeconomie beschrijven, inclusief
kansen, risico's, sociale en ethische implicaties. & Gemiddeld & K &
AI-Impliciet \\
LO2.3.18 & Het concept en de functies van maatschappelijke monitoring
beschrijven. & Gemiddeld & K & - \\
LO2.3.19 & Het concept en de functies van e-overheid beschrijven. &
Gemiddeld & K & - \\
LO2.3.20 & Beschrijven hoe digitale technologieën zoals sociale
mediaplatforms bepaalde aspecten van de basisdemocratie kunnen
beïnvloeden (bijvoorbeeld door beïnvloeding van het verkiezingsproces).
& Gemiddeld & K & AI-Impliciet \\
LO2.3.21 & Autonoom en effectief interageren met digitale platforms en
diensten. & Gemiddeld & V & AI-Impliciet \\
LO2.3.22 & Kansen, risico's, sociale en ethische implicaties van de
platformeconomie beoordelen. & Gemiddeld & V & AI-Impliciet \\
LO2.3.23 & Prioriteit geven aan het voortdurend verkennen van manieren
waarop digitale technologieën empowerment of maatschappelijke
participatie kunnen ondersteunen. & Geavanceerd & H & AI-Impliciet \\
LO2.3.24 & Deelnemen aan discussies over de ethische, politieke en
sociale implicaties van digitale technologieën. & Geavanceerd & H &
AI-Impliciet \\
LO2.3.25 & Onderscheid maken tussen hoogrisico en verboden AI-systemen
(volgens wetgeving). & Geavanceerd & K & AI-Expliciet \\
LO2.3.26 & Mogelijke maatschappelijke, politieke of economische effecten
van verboden en hoogrisico AI-systemen beschrijven. & Geavanceerd & K &
AI-Expliciet \\
LO2.3.27 & Het potentieel van digitale technologieën voor inclusie,
uitsluiting en maatschappelijke interventie in een gegeven context
beoordelen. & Geavanceerd & V & AI-Impliciet \\
LO2.3.28 & Een reeks manieren beoordelen waarop digitale technologieën
zoals sociale mediaplatforms democratische processen kunnen beïnvloeden.
& Geavanceerd & V & AI-Impliciet \\
LO2.3.29 & Anderen helpen bij het identificeren van kansen en het
participeren in digitale omgevingen voor empowerment en participatie
(van zichzelf of de gemeenschap). & Geavanceerd & V & - \\
LO2.3.30 & Anderen ondersteunen bij het zich informeren over en het
uitoefenen van hun rechten onder digitale wetgeving. & Geavanceerd & V &
AI-Impliciet \\
LO2.3.31 & Op de hoogte blijven van vrijheden, rechten en
verantwoordelijkheden van individuen bij evoluerende digitale
technologieën en wetsontwikkelingen. & Zeer geavanceerd & H &
AI-Impliciet \\
LO2.3.32 & Meerdere effecten van digitale technologieën op de
samenleving, politieke processen of de economie evalueren vanuit
verschillende perspectieven. & Zeer geavanceerd & V & AI-Impliciet \\
LO2.3.33 & Anderen helpen de belangrijkste bepalingen van digitale
wetgeving te begrijpen, gegeven een specifieke context. & Zeer
geavanceerd & V & AI-Impliciet \\
LO2.3.34 & Leiding geven aan of ontwerpen van digitale
burgerschapsinitiatieven, bijvoorbeeld om participatie, inclusie of
empowerment te bevorderen. & Zeer geavanceerd & V & AI-Impliciet \\
\end{longtable}

\subsubsection{2.4 Samenwerken via digitale
technologieën}\label{samenwerken-via-digitale-technologieuxebn}

\begin{longtable}[]{@{}
  >{\centering\arraybackslash}p{(\linewidth - 8\tabcolsep) * \real{0.1053}}
  >{\centering\arraybackslash}p{(\linewidth - 8\tabcolsep) * \real{0.3947}}
  >{\centering\arraybackslash}p{(\linewidth - 8\tabcolsep) * \real{0.2105}}
  >{\centering\arraybackslash}p{(\linewidth - 8\tabcolsep) * \real{0.1842}}
  >{\centering\arraybackslash}p{(\linewidth - 8\tabcolsep) * \real{0.1053}}@{}}
\toprule\noalign{}
\begin{minipage}[b]{\linewidth}\centering
ID
\end{minipage} & \begin{minipage}[b]{\linewidth}\centering
Leerresultaat
\end{minipage} & \begin{minipage}[b]{\linewidth}\centering
Niveau
\end{minipage} & \begin{minipage}[b]{\linewidth}\centering
K/V/H
\end{minipage} & \begin{minipage}[b]{\linewidth}\centering
AI
\end{minipage} \\
\midrule\noalign{}
\endhead
\bottomrule\noalign{}
\endlastfoot
LO2.4.01 & Het belang inzien van effectieve communicatieve vaardigheden
voor succesvolle samenwerking in digitale omgevingen. & Basis & H & - \\
LO2.4.02 & Belangrijkste voordelen en beperkingen herkennen van digitale
samenwerkingshulpmiddelen. & Basis & K & AI-Impliciet \\
LO2.4.03 & De aanwezigheid van AI-systemen in digitale
samenwerkingshulpmiddelen herkennen. & Basis & K & AI-Expliciet \\
LO2.4.04 & Deelnemen aan samenwerkingsgroepen via digitale
samenwerkingshulpmiddelen. & Basis & V & AI-Impliciet \\
LO2.4.05 & Rekening houden met verschillende perspectieven om te helpen
een gemeenschappelijk doel te bereiken in digitale omgevingen. &
Gemiddeld & H & - \\
LO2.4.06 & Belangrijkste kenmerken en functies van verschillende
digitale samenwerkingshulpmiddelen identificeren. & Gemiddeld & K &
AI-Impliciet \\
LO2.4.07 & De functies, voordelen en beperkingen herkennen van
AI-systeemfunctionaliteiten in sommige digitale
samenwerkingshulpmiddelen. & Gemiddeld & K & AI-Expliciet \\
LO2.4.08 & Voorbeelden identificeren van ethische, verantwoordelijke en
effectieve samenwerking tussen mens en AI. & Gemiddeld & K &
AI-Expliciet \\
LO2.4.09 & Digitale samenwerkingshulpmiddelen selecteren die passen bij
de samenwerkingsdoelen. & Gemiddeld & V & AI-Impliciet \\
LO2.4.10 & Eenvoudige samenwerkingstaken creëren en beheren met behulp
van digitale samenwerkingshulpmiddelen. & Gemiddeld & V &
AI-Impliciet \\
LO2.4.11 & Effectief bijdragen aan eenvoudige samenwerkingstaken in
digitale omgevingen. & Gemiddeld & V & AI-Impliciet \\
LO2.4.12 & Prioriteit geven aan een goede afstemming tussen het gebruik
van digitale samenwerkingshulpmiddelen en de voorkeuren van de betrokken
individuen. & Geavanceerd & H & AI-Impliciet \\
LO2.4.13 & Zorgen voor passend en ethisch gebruik van Digitale
technologieën, waaronder AI-systemen voor samenwerkingsactiviteiten. &
Geavanceerd & H & AI-Expliciet \\
LO2.4.14 & Verschillende digitale samenwerkingshulpmiddelen gebruiken en
combineren die voldoen aan de behoeften van projecten, taken en groepen.
& Geavanceerd & V & AI-Impliciet \\
LO2.4.15 & Anderen helpen bij het ontwikkelen van hun vaardigheden om
samen te werken in digitale omgevingen. & Geavanceerd & V &
AI-Impliciet \\
LO2.4.16 & Ethische en praktische aspecten van
mens-AI-samenwerkingstechnieken voor een bepaald doel beoordelen. &
Geavanceerd & V & AI-Expliciet \\
LO2.4.17 & Leiding geven aan samenwerking in digitale omgevingen. &
Geavanceerd & V & AI-Impliciet \\
LO2.4.18 & Op de hoogte blijven van ontwikkelingen in
samenwerkingspraktijken in digitale omgevingen. & Zeer geavanceerd & H &
AI-Impliciet \\
LO2.4.19 & Proportioneel, ethisch en effectief gebruik van digitale
technologieën, inclusief AI-systemen, in samenwerkingen bevorderen en
ondersteunen. & Zeer geavanceerd & H & AI-Expliciet \\
LO2.4.20 & Complexe of gespecialiseerde strategieën voor samenwerking in
digitale omgevingen ontwerpen. & Zeer geavanceerd & V & AI-Impliciet \\
LO2.4.21 & Anderen helpen bij het ontwikkelen van vaardigheden om
leiding te geven aan samenwerking in digitale omgevingen. & Zeer
geavanceerd & V & AI-Impliciet \\
LO2.4.22 & Leiden van of bijdragen aan verbeteringen in of nieuwe
oplossingen voor mens-AI-samenwerking. & Zeer geavanceerd & V &
AI-Expliciet \\
\end{longtable}

\subsubsection{2.5 Digitaal gedrag}\label{digitaal-gedrag}

\begin{longtable}[]{@{}
  >{\centering\arraybackslash}p{(\linewidth - 8\tabcolsep) * \real{0.1053}}
  >{\centering\arraybackslash}p{(\linewidth - 8\tabcolsep) * \real{0.3947}}
  >{\centering\arraybackslash}p{(\linewidth - 8\tabcolsep) * \real{0.2105}}
  >{\centering\arraybackslash}p{(\linewidth - 8\tabcolsep) * \real{0.1842}}
  >{\centering\arraybackslash}p{(\linewidth - 8\tabcolsep) * \real{0.1053}}@{}}
\toprule\noalign{}
\begin{minipage}[b]{\linewidth}\centering
ID
\end{minipage} & \begin{minipage}[b]{\linewidth}\centering
Leerresultaat
\end{minipage} & \begin{minipage}[b]{\linewidth}\centering
Niveau
\end{minipage} & \begin{minipage}[b]{\linewidth}\centering
K/V/H
\end{minipage} & \begin{minipage}[b]{\linewidth}\centering
AI
\end{minipage} \\
\midrule\noalign{}
\endhead
\bottomrule\noalign{}
\endlastfoot
LO2.5.01 & Het belang inzien van het geven van ruimte aan de meningen
van anderen in digitale omgevingen. & Basis & H & - \\
LO2.5.02 & Verschillen identificeren in verbaal en non-verbaal gedrag in
digitale en niet-digitale omgevingen. & Basis & K & - \\
LO2.5.03 & Herkennen dat er culturele en contextuele verschillen zijn in
verbale en non-verbale communicatie in digitale omgevingen. & Basis & K
& - \\
LO2.5.04 & Herkennen dat bepaald gedrag in digitale omgevingen mogelijk
niet acceptabel is voor anderen en/of juridische gevolgen kan hebben. &
Basis & K & AI-Impliciet \\
LO2.5.05 & Passende toon en visuele expressie zoals emoji gebruiken in
formele en informele digitale omgevingen. & Basis & V & - \\
LO2.5.06 & Prioriteit geven aan digitaal gedrag dat inclusie en een
positieve digitale reputatie van zichzelf en anderen ondersteunt. &
Gemiddeld & H & - \\
LO2.5.07 & De belangrijkste rechten en verantwoordelijkheden van
kinderen en volwassenen met betrekking tot digitaal gedrag
identificeren. & Gemiddeld & K & AI-Impliciet \\
LO2.5.08 & De relatie tussen digitaal gedrag en digitale reputatie
beschrijven. & Gemiddeld & K & - \\
LO2.5.09 & De belangrijkste rechten en verantwoordelijkheden van
kinderen en volwassenen met betrekking tot digitaal gedrag
identificeren. & Gemiddeld & K & AI-Impliciet \\
LO2.5.10 & Een respectvolle en inclusieve toon en visuele expressie
zoals emoji gebruiken in formele en informele digitale omgevingen. &
Gemiddeld & V & - \\
LO2.5.11 & Inclusief en respectvol gedrag in digitale omgevingen
bevorderen en ondersteunen. & Geavanceerd & H & - \\
LO2.5.12 & Onderscheid maken tussen ethisch, legaal en illegaal gedrag
in digitale omgevingen, erkennend dat deze onderscheiden complex kunnen
zijn. & Geavanceerd & K & AI-Impliciet \\
LO2.5.13 & Soorten misbruik identificeren die in digitale omgevingen
kunnen voorkomen, inclusief getroffen groepen en mogelijke effecten. &
Geavanceerd & K & AI-Impliciet \\
LO2.5.14 & Manieren beschrijven waarop misbruik in digitale omgevingen
kan worden gemeld en aangepakt. & Geavanceerd & K & AI-Impliciet \\
LO2.5.15 & Met effectieve en respectvolle communicatie en gedrag
reageren op moeilijke of complexe situaties in digitale omgevingen. &
Geavanceerd & V & - \\
LO2.5.16 & Ethische en legale/illegale aspecten van gedrag in digitale
omgevingen in een specifieke context beoordelen. & Geavanceerd & V &
AI-Impliciet \\
LO2.5.17 & Patronen en mogelijke effecten van misbruik van specifieke
groepen in digitale omgevingen analyseren. & Geavanceerd & V &
AI-Impliciet \\
LO2.5.18 & Anderen ondersteunen bij het ontwikkelen van hun vaardigheden
voor inclusief en respectvol gedrag in digitale omgevingen. &
Geavanceerd & V & - \\
LO2.5.19 & Op de hoogte blijven van ontwikkelingen in beleid en
wetgeving met betrekking tot gedrag in digitale omgevingen. & Zeer
geavanceerd & H & AI-Impliciet \\
LO2.5.20 & Anderen helpen belangrijke rechten en verantwoordelijkheden
te begrijpen onder beleid en wetgeving met betrekking tot digitaal
gedrag in een gegeven context. & Zeer geavanceerd & V & AI-Impliciet \\
LO2.5.21 & Leiden van of bijdragen aan beleid of initiatieven op het
gebied van digitaal gedrag. & Zeer geavanceerd & V & AI-Impliciet \\
\end{longtable}

\subsubsection{2.6 Digitale identiteit
beheren}\label{digitale-identiteit-beheren}

\begin{longtable}[]{@{}
  >{\centering\arraybackslash}p{(\linewidth - 8\tabcolsep) * \real{0.1053}}
  >{\centering\arraybackslash}p{(\linewidth - 8\tabcolsep) * \real{0.3947}}
  >{\centering\arraybackslash}p{(\linewidth - 8\tabcolsep) * \real{0.2105}}
  >{\centering\arraybackslash}p{(\linewidth - 8\tabcolsep) * \real{0.1842}}
  >{\centering\arraybackslash}p{(\linewidth - 8\tabcolsep) * \real{0.1053}}@{}}
\toprule\noalign{}
\begin{minipage}[b]{\linewidth}\centering
ID
\end{minipage} & \begin{minipage}[b]{\linewidth}\centering
Leerresultaat
\end{minipage} & \begin{minipage}[b]{\linewidth}\centering
Niveau
\end{minipage} & \begin{minipage}[b]{\linewidth}\centering
K/V/H
\end{minipage} & \begin{minipage}[b]{\linewidth}\centering
AI
\end{minipage} \\
\midrule\noalign{}
\endhead
\bottomrule\noalign{}
\endlastfoot
LO2.6.01 & De voordelen inzien van het implementeren van maatregelen om
de eigen digitale identiteit te helpen beheren. & Basis & H & - \\
LO2.6.02 & Kenmerken van fysieke en digitale identiteiten herkennen. &
Basis & K & - \\
LO2.6.03 & Aspecten van fysieke identiteit identificeren die gekoppeld
kunnen worden aan digitale identiteit. & Basis & K & - \\
LO2.6.04 & Digitale identiteit herkennen als zowel een methode om een
individu te authenticeren (verifiëren) als de data gegenereerd door de
online activiteiten van een individu. & Basis & K & AI-Impliciet \\
LO2.6.05 & Veelvoorkomende vormen en toepassingen van digitale
identiteit identificeren. & Basis & K & AI-Impliciet \\
LO2.6.06 & Herkennen dat informatie op het internet in de loop van de
tijd kan blijven bestaan. & Basis & K & AI-Impliciet \\
LO2.6.07 & Het concept en de componenten van een digitale voetafdruk
herkennen. & Basis & K & AI-Impliciet \\
LO2.6.08 & Herkennen dat wetten voor de bescherming van digitale
identiteit de data en privacy van individuen beschermen. & Basis & K &
AI-Impliciet \\
LO2.6.09 & Breng eenvoudige maatregelen in kaart, zoals het beperken van
tracking en het verwijderen van de browsergeschiedenis, om de digitale
identiteit te beheren. & Basis & K & - \\
LO2.6.10 & Neem eenvoudige maatregelen, zoals het beperken van tracking
en het verwijderen van de browsergeschiedenis, om de digitale identiteit
te beheren. & Basis & V & - \\
LO2.6.11 & Het belang inzien van de eigen rol en rechten bij het beheer
van de digitale identiteit. & Gemiddeld & H & AI-Impliciet \\
LO2.6.12 & De relaties beschrijven tussen digitale voetafdruk, digitale
reputatie en digitale identiteit. & Gemiddeld & K & AI-Impliciet \\
LO2.6.13 & Voorbeelden identificeren van actief en passief gegenereerde
informatie in relatie tot digitale identiteit. & Gemiddeld & K &
AI-Impliciet \\
LO2.6.14 & Manieren beschrijven waarop de reikwijdte van de eigen
digitale identiteit kan worden geanalyseerd. & Gemiddeld & K &
AI-Impliciet \\
LO2.6.15 & Kenmerken en functies identificeren die worden gebruikt om de
digitale identiteit te beheren, zoals instellingen op apparaten en apps,
online accounts, activiteitstracking en sociale mediaplatforms. &
Gemiddeld & K & AI-Impliciet \\
LO2.6.16 & Informatie over de reikwijdte van de eigen digitale
identiteit gebruiken om acties op het gebied van identiteitsbeheer te
sturen. & Gemiddeld & V & AI-Impliciet \\
LO2.6.17 & Instellingen aanpassen op apparaten en apps, online accounts
en activiteitstracking om de eigen digitale identiteit te helpen
beheren. & Gemiddeld & V & AI-Impliciet \\
LO2.6.18 & Een of meer digitale identiteiten cureren en beheren met
behulp van verschillende kenmerken en functionaliteiten op digitale
platform(s) of diensten. & Gemiddeld & V & AI-Impliciet \\
LO2.6.19 & Prioriteit geven aan de doorlopende beoordeling van de eigen
digitale identiteit. & Geavanceerd & H & AI-Impliciet \\
LO2.6.20 & De relatie herkennen tussen digitale technologische
ontwikkelingen en het beheer van de digitale identiteit. & Geavanceerd &
K & AI-Impliciet \\
LO2.6.21 & Manieren beschrijven om rechten uit te oefenen op het gebied
van digitale identiteit. & Geavanceerd & K & AI-Impliciet \\
LO2.6.22 & Manieren beschrijven waarop AI-systemen worden gebruikt bij
het beheer van digitale identiteit. & Geavanceerd & K & AI-Expliciet \\
LO2.6.23 & Verschillende processen implementeren om de digitale
identiteit over een reeks digitale omgevingen te beheren. & Geavanceerd
& V & AI-Impliciet \\
LO2.6.24 & Voordelen, sociale en ethische implicaties van het gebruik
van AI-systemen bij het beheer van digitale identiteit beoordelen. &
Geavanceerd & V & AI-Expliciet \\
LO2.6.25 & Digitale identiteiten cureren en beheren voor persoonlijke,
professionele en/of organisatorische doeleinden over verschillende
platforms en diensten. & Geavanceerd & V & AI-Impliciet \\
LO2.6.26 & Anderen helpen bij het basisbeheer van hun digitale
identiteit. & Geavanceerd & V & AI-Impliciet \\
LO2.6.27 & Op de hoogte blijven van ontwikkelingen in digitale
technologieën in relatie tot het beheer van de digitale identiteit. &
Zeer geavanceerd & H & AI-Impliciet \\
LO2.6.28 & Anderen ondersteunen bij het verdiepen van hun vaardigheden
in het beheer en de curatie van digitale identiteiten. & Zeer
geavanceerd & V & AI-Impliciet \\
LO2.6.29 & Anderen adviseren over complexe aspecten van het beheer van
digitale identiteit en bijbehorende rechten. & Zeer geavanceerd & V &
AI-Impliciet \\
\end{longtable}

\subsubsection{3.1 Digitale inhoud
ontwikkelen}\label{digitale-inhoud-ontwikkelen}

\begin{longtable}[]{@{}
  >{\centering\arraybackslash}p{(\linewidth - 8\tabcolsep) * \real{0.1053}}
  >{\centering\arraybackslash}p{(\linewidth - 8\tabcolsep) * \real{0.3947}}
  >{\centering\arraybackslash}p{(\linewidth - 8\tabcolsep) * \real{0.2105}}
  >{\centering\arraybackslash}p{(\linewidth - 8\tabcolsep) * \real{0.1842}}
  >{\centering\arraybackslash}p{(\linewidth - 8\tabcolsep) * \real{0.1053}}@{}}
\toprule\noalign{}
\begin{minipage}[b]{\linewidth}\centering
ID
\end{minipage} & \begin{minipage}[b]{\linewidth}\centering
Leerresultaat
\end{minipage} & \begin{minipage}[b]{\linewidth}\centering
Niveau
\end{minipage} & \begin{minipage}[b]{\linewidth}\centering
K/V/H
\end{minipage} & \begin{minipage}[b]{\linewidth}\centering
AI
\end{minipage} \\
\midrule\noalign{}
\endhead
\bottomrule\noalign{}
\endlastfoot
LO3.1.01 & De voordelen inzien van het verkennen van verschillende
hulpmiddelen voor het creëren van digitale inhoud ter ondersteuning van
contentcreatiedoelen. & Basis & H & AI-Impliciet \\
LO3.1.02 & Het belang inzien van digitale inhoud die toegankelijk en
inclusief is. & Basis & H & - \\
LO3.1.03 & Veelvoorkomende soorten digitale inhoud en hun bijbehorende
bestandsformaten identificeren. & Basis & K & - \\
LO3.1.04 & Veelvoorkomende operationele functies identificeren binnen
hulpmiddelen voor het creëren van digitale inhoud. & Basis & K &
AI-Impliciet \\
LO3.1.05 & Onderscheid maken tussen toegankelijke digitale inhoud en
inclusieve digitale inhoud. & Basis & K & - \\
LO3.1.06 & Herkennen dat AI-systemen weliswaar digitale inhoud kunnen
genereren en integreren, maar dat mensen essentieel zijn om ethische,
verantwoordelijke en context-passende outputs te waarborgen. & Basis & K
& AI-Expliciet \\
LO3.1.07 & Herkennen dat generatieve AI een specifiek type AI is en een
van de verschillende digitale technologieën die gebruikt kunnen worden
om contentcreatie te ondersteunen. & Basis & K & AI-Expliciet \\
LO3.1.08 & Basiskenmerken van hulpmiddelen voor het creëren van digitale
inhoud gebruiken om digitale inhoud te creëren en te bewerken (tekst,
afbeelding, video en/of audio). & Basis & V & AI-Impliciet \\
LO3.1.09 & Doelgericht kenmerken en functies van hulpmiddelen voor het
creëren van digitale inhoud verkennen om contentcreatievaardigheden te
verdiepen. & Gemiddeld & H & AI-Impliciet \\
LO3.1.10 & Voordelen, beperkingen en ethische overwegingen beschrijven
bij het gebruik van digitale technologieën zoals AI-systemen voor
contentcreatie. & Gemiddeld & K & AI-Expliciet \\
LO3.1.11 & Bepaal strategieën zoals sjablonen of een passende volgorde
van stappen die zorgen voor een efficiënte creatie van digitale inhoud.
& Gemiddeld & K & AI-Impliciet \\
LO3.1.12 & Verschillende hulpmiddelen voor het creëren van digitale
inhoud gebruiken om digitale inhoud te creëren en te bewerken (tekst,
afbeelding, video en/of audio). & Gemiddeld & V & AI-Impliciet \\
LO3.1.13 & De behoeften op het gebied van inclusiviteit en/of
toegankelijkheid beoordelen van het publiek waarvoor digitale inhoud
wordt gecreëerd. & Gemiddeld & V & - \\
LO3.1.14 & Digitale inhoud bewerken om de toegankelijkheid te verbeteren
en te voldoen aan de behoeften van het publiek. & Gemiddeld & V &
AI-Impliciet \\
LO3.1.15 & Pas strategieën toe, zoals sjablonen of een juiste volgorde
van stappen die de efficiënte creatie van digitale inhoud mogelijk
maken. & Gemiddeld & V & AI-Impliciet \\
LO3.1.16 & Doelgericht, selectief en ethisch interageren met AI-systemen
ter ondersteuning van de creatie van digitale inhoud. & Gemiddeld & V &
AI-Expliciet \\
LO3.1.17 & Het belang inzien van het beoordelen van mogelijkheden,
beperkingen en ethische aspecten van hulpmiddelen voor het creëren van
digitale inhoud om een passende selectie en gebruik te waarborgen. &
Geavanceerd & H & AI-Impliciet \\
LO3.1.18 & Hulpmiddelen en methoden voor het creëren van digitale inhoud
selecteren en combineren om te voldoen aan complexe contentcreatietaken
en publiekseisen. & Geavanceerd & V & AI-Impliciet \\
LO3.1.19 & Een verscheidenheid aan complexe of gespecialiseerde digitale
inhoud creëren en bewerken, passend afgestemd op doelen en publiek. &
Geavanceerd & V & AI-Impliciet \\
LO3.1.20 & Anderen ondersteunen bij het ontwikkelen van hun vaardigheden
in de creatie van digitale inhoud door middel van ethische en
verantwoordelijke benaderingen. & Geavanceerd & V & AI-Impliciet \\
LO3.1.21 & Toegankelijkheid en inclusiviteit in initiatieven voor de
creatie van digitale inhoud bevorderen en ondersteunen. & Zeer
geavanceerd & H & AI-Impliciet \\
LO3.1.22 & Het selectieve en ethische gebruik van AI-systemen bij
contentcreatie bevorderen en ondersteunen. & Zeer geavanceerd & H &
AI-Expliciet \\
LO3.1.23 & Anderen helpen bij het ontwikkelen van geavanceerde
vaardigheden in de creatie van digitale inhoud. & Zeer geavanceerd & V &
AI-Impliciet \\
LO3.1.24 & Leiding geven aan of bijdragen aan complexe of
gespecialiseerde initiatieven voor de creatie van digitale inhoud. &
Zeer geavanceerd & V & AI-Impliciet \\
LO3.1.25 & Leiden van of bijdragen aan verbeteringen in of nieuwe
oplossingen voor complexe of gespecialiseerde digitale inhoud. & Zeer
geavanceerd & V & AI-Impliciet \\
\end{longtable}

\subsubsection{3.2 Integreren en herwerken van digitale
inhoud}\label{integreren-en-herwerken-van-digitale-inhoud}

\begin{longtable}[]{@{}
  >{\centering\arraybackslash}p{(\linewidth - 8\tabcolsep) * \real{0.1053}}
  >{\centering\arraybackslash}p{(\linewidth - 8\tabcolsep) * \real{0.3947}}
  >{\centering\arraybackslash}p{(\linewidth - 8\tabcolsep) * \real{0.2105}}
  >{\centering\arraybackslash}p{(\linewidth - 8\tabcolsep) * \real{0.1842}}
  >{\centering\arraybackslash}p{(\linewidth - 8\tabcolsep) * \real{0.1053}}@{}}
\toprule\noalign{}
\begin{minipage}[b]{\linewidth}\centering
ID
\end{minipage} & \begin{minipage}[b]{\linewidth}\centering
Leerresultaat
\end{minipage} & \begin{minipage}[b]{\linewidth}\centering
Niveau
\end{minipage} & \begin{minipage}[b]{\linewidth}\centering
K/V/H
\end{minipage} & \begin{minipage}[b]{\linewidth}\centering
AI
\end{minipage} \\
\midrule\noalign{}
\endhead
\bottomrule\noalign{}
\endlastfoot
LO3.2.01 & Het belang inzien van ethische en transparante praktijken bij
het hergebruiken of uitwerken van bestaande digitale inhoud. & Basis & H
& AI-Impliciet \\
LO3.2.02 & De voordelen inzien van het verkennen van hulpmiddelen en
technieken voor de integratie en uitwerking van digitale inhoud. & Basis
& H & AI-Impliciet \\
LO3.2.03 & Het concept van bronvermelding herkennen bij het hergebruiken
van bestaande digitale inhoud. & Basis & K & AI-Impliciet \\
LO3.2.04 & Het concept van transparant gebruik van AI-systemen (met name
generatieve AI) herkennen bij de integratie en herbewerking van digitale
inhoud. & Basis & K & AI-Expliciet \\
LO3.2.05 & Onderscheid maken tussen bewerkbare en niet-bewerkbare
digitale inhoud. & Basis & K & - \\
LO3.2.06 & Belangrijkste functies van hulpmiddelen voor contentcreatie
identificeren voor het bewerken en integreren van digitale inhoud
(tekst, afbeelding, audio, video). & Basis & K & AI-Impliciet \\
LO3.2.07 & Wijzigingen aanbrengen in bestaande digitale inhoud met
behulp van basisbewerkings-, formatterings- en integratiefuncties. &
Basis & V & AI-Impliciet \\
LO3.2.08 & Doelgericht verschillende manieren verkennen om digitale
inhoud te integreren en te herbewerken om vaardigheden op dit gebied te
verdiepen. & Gemiddeld & H & AI-Impliciet \\
LO3.2.09 & Structuur-, formaat- en publiekseisen identificeren van een
taak voor integratie of herbewerking van digitale inhoud. & Gemiddeld &
K & - \\
LO3.2.10 & Manieren beschrijven waarop bronvermelding kan worden
toegepast voor digitale inhoud die wordt hergebruikt. & Gemiddeld & K &
- \\
LO3.2.11 & Ethische en transparante praktijken beschrijven bij het
gebruik van AI-systemen (met name generatieve AI) bij de integratie en
uitwerking van digitale inhoud. & Gemiddeld & K & AI-Expliciet \\
LO3.2.12 & Digitale inhoud aanpassen of integreren om te voldoen aan
formaat-, structuur- en publiekseisen. & Gemiddeld & V & AI-Impliciet \\
LO3.2.13 & Digitale tekstuele, numerieke of visuele representaties
wijzigen of transformeren om de betekenis van data en informatie
effectief en nauwkeurig over te brengen. & Gemiddeld & V &
AI-Impliciet \\
LO3.2.14 & Digitale technologieën op een selectieve, ethische en
transparante manier gebruiken om verbeteringen of integraties aan te
brengen in bestaande digitale inhoud. & Gemiddeld & V & AI-Impliciet \\
LO3.2.15 & Prioriteit geven aan transparante en ethische praktijken bij
taken voor integratie en herbewerking van digitale inhoud. & Geavanceerd
& H & AI-Impliciet \\
LO3.2.16 & Een reeks methoden beschrijven voor complexe integratie en
herbewerking van digitale inhoud. & Geavanceerd & K & AI-Impliciet \\
LO3.2.17 & Passend en onpassend gebruik van AI-systemen beschrijven om
de integratie of herbewerking van digitale inhoud bij complexe taken te
verbeteren. & Geavanceerd & K & AI-Expliciet \\
LO3.2.18 & Verschillende soorten digitale inhoud aanpassen of integreren
om te voldoen aan complexe formaat-, structuur- en publiekseisen. &
Geavanceerd & V & AI-Impliciet \\
LO3.2.19 & Digitale technologieën op een selectieve, ethische en
transparante manier inzetten voor complexe integratie of herbewerking
van digitale inhoud. & Geavanceerd & V & AI-Impliciet \\
LO3.2.20 & Anderen ondersteunen bij het ontwikkelen van hun vaardigheden
in de integratie en herbewerking van digitale inhoud. & Geavanceerd & V
& AI-Impliciet \\
LO3.2.21 & Ethische en transparante praktijken bij de integratie en
herbewerking van digitale inhoud bevorderen en ondersteunen. & Zeer
geavanceerd & H & AI-Impliciet \\
LO3.2.22 & Op de hoogte blijven van technologische ontwikkelingen op het
gebied van integratie en herbewerking van digitale inhoud, inclusief hun
technische en ethische implicaties. & Zeer geavanceerd & H &
AI-Impliciet \\
LO3.2.23 & Geavanceerde ontwerp- en datavisualisatietechnieken evalueren
en toepassen op complexe of gespecialiseerde integratie en herbewerking
van digitale inhoud. & Zeer geavanceerd & V & AI-Impliciet \\
LO3.2.24 & Anderen helpen bij complexe taken voor integratie en
herbewerking van digitale inhoud. & Zeer geavanceerd & V &
AI-Impliciet \\
LO3.2.25 & Leiding geven aan of bijdragen aan complexe initiatieven voor
integratie of herbewerking van digitale inhoud. & Zeer geavanceerd & V &
AI-Impliciet \\
LO3.2.26 & Leiden van of bijdragen aan verbeteringen in of nieuwe
oplossingen voor complexe integratie of herbewerking van digitale
inhoud. & Zeer geavanceerd & V & AI-Impliciet \\
\end{longtable}

\subsubsection{3.3 Auteursrecht en
licenties}\label{auteursrecht-en-licenties}

\begin{longtable}[]{@{}
  >{\centering\arraybackslash}p{(\linewidth - 8\tabcolsep) * \real{0.1053}}
  >{\centering\arraybackslash}p{(\linewidth - 8\tabcolsep) * \real{0.3947}}
  >{\centering\arraybackslash}p{(\linewidth - 8\tabcolsep) * \real{0.2105}}
  >{\centering\arraybackslash}p{(\linewidth - 8\tabcolsep) * \real{0.1842}}
  >{\centering\arraybackslash}p{(\linewidth - 8\tabcolsep) * \real{0.1053}}@{}}
\toprule\noalign{}
\begin{minipage}[b]{\linewidth}\centering
ID
\end{minipage} & \begin{minipage}[b]{\linewidth}\centering
Leerresultaat
\end{minipage} & \begin{minipage}[b]{\linewidth}\centering
Niveau
\end{minipage} & \begin{minipage}[b]{\linewidth}\centering
K/V/H
\end{minipage} & \begin{minipage}[b]{\linewidth}\centering
AI
\end{minipage} \\
\midrule\noalign{}
\endhead
\bottomrule\noalign{}
\endlastfoot
LO3.3.01 & Herkennen dat het internet geen volledig vrije ruimte is: de
data van individuen kan te gelde worden gemaakt en individuen hebben
mogelijk toestemming nodig om gevonden inhoud te gebruiken. & Basis & K
& AI-Impliciet \\
LO3.3.02 & In algemene termen de concepten van auteursrecht en licenties
herkennen in digitale contexten. & Basis & K & - \\
LO3.3.03 & Herkennen dat op de originele digitale inhoud van een
individu automatisch auteursrecht rust. & Basis & K & - \\
LO3.3.04 & Herkennen dat er verschillende soorten auteursrechten en
licenties van toepassing zijn op digitale inhoud, inclusief door AI
gegenereerde inhoud, en dat deze bepalen hoe inhoud kan worden gebruikt
en gedeeld. & Basis & K & AI-Expliciet \\
LO3.3.05 & Herkennen dat door AI gegenereerde inhoud als zodanig moet
worden gelabeld om anderen te helpen de oorsprong en mogelijkheden voor
verder gebruik te begrijpen. & Basis & K & AI-Expliciet \\
LO3.3.06 & Digitale inhoud identificeren die gratis kan worden gebruikt.
& Basis & K & - \\
LO3.3.07 & Digitale inhoud gebruiken en delen in overeenstemming met de
basisregels van juridische en ethische richtlijnen. & Basis & V &
AI-Impliciet \\
LO3.3.08 & Prioriteit geven aan een voorzichtige benadering van digitale
inhoud (controleren voor gebruik of delen). & Gemiddeld & H &
AI-Impliciet \\
LO3.3.09 & De complexe aard van auteursrecht en licenties van digitale
inhoud inzien. & Gemiddeld & H & AI-Impliciet \\
LO3.3.10 & Het concept van intellectueel eigendom definiëren, met
voorbeelden uit digitale contexten. & Gemiddeld & K & AI-Impliciet \\
LO3.3.11 & Onderscheid maken, met voorbeelden uit digitale contexten,
tussen auteursrecht, handelsmerk, modelrecht en octrooi. & Gemiddeld & K
& AI-Impliciet \\
LO3.3.12 & Veelvoorkomende soorten en doelen van licenties in digitale
contexten identificeren, inclusief Creative Commons. & Gemiddeld & K &
- \\
LO3.3.13 & Voorbeelden van piraterij en plagiaat in digitale contexten
identificeren. & Gemiddeld & K & AI-Impliciet \\
LO3.3.14 & Voorbeelden identificeren van waar auteursrecht wel en niet
van toepassing is in digitale contexten. & Gemiddeld & K &
AI-Impliciet \\
LO3.3.15 & Juridische, ethische en commerciële gevolgen beschrijven van
inbreuken op intellectueel eigendom in digitale contexten, inclusief
piraterij en plagiaat. & Gemiddeld & K & AI-Impliciet \\
LO3.3.16 & Voorbeelden identificeren van juridische en ethische
uitdagingen met betrekking tot auteursrecht bij het trainen van
AI-modellen. & Gemiddeld & K & AI-Expliciet \\
LO3.3.17 & Juridische en ethische richtlijnen op een passende manier
toepassen bij het gebruiken en delen van digitale inhoud. & Gemiddeld &
V & AI-Impliciet \\
LO3.3.18 & Prioriteit geven aan het meewegen van ethische en juridische
aspecten, zoals transparantie en auteursrecht, bij het gebruiken en
delen van digitale inhoud. & Geavanceerd & H & AI-Impliciet \\
LO3.3.19 & Belangrijkste kenmerken van de huidige wetgeving met
betrekking tot digitaal auteursrecht en licenties beschrijven. &
Geavanceerd & K & AI-Impliciet \\
LO3.3.20 & Voorbeelden beschrijven van waar auteursrecht wel en niet van
toepassing is in digitale contexten. & Geavanceerd & K & AI-Impliciet \\
LO3.3.21 & Verschillen identificeren in hoe ethische en
auteursrechtelijke kwesties van toepassing zijn op trainingsdata voor
AI-systemen en door AI gegenereerde inhoud (output). & Geavanceerd & K &
AI-Expliciet \\
LO3.3.22 & Juridische en ethische richtlijnen voor het gebruiken en
delen van digitale inhoud in complexe contexten beoordelen en correct
toepassen, inclusief verschillende softwarelicentiemodellen en de
vernieuwingsvereisten van licenties. & Geavanceerd & V & AI-Impliciet \\
LO3.3.23 & Anderen helpen bij het gebruiken en delen van digitale inhoud
in overeenstemming met de juridische en ethische richtlijnen. &
Geavanceerd & V & AI-Impliciet \\
LO3.3.24 & Op de hoogte blijven van ontwikkelingen in auteursrecht- en
licentieregelgeving in digitale contexten. & Zeer geavanceerd & H &
AI-Impliciet \\
LO3.3.25 & Bewustzijn en begrip van juridische en ethische auteursrecht-
en licentiepraktijken in digitale contexten bevorderen en ondersteunen.
& Zeer geavanceerd & H & AI-Impliciet \\
LO3.3.26 & Geavanceerde kennis van intellectuele eigendomsrechten,
auteursrecht en licentieconcepten in digitale contexten toepassen ter
ondersteuning van besluitvorming. & Zeer geavanceerd & V &
AI-Impliciet \\
LO3.3.27 & Leiding geven aan of bijdragen aan beleid of richtlijnen over
auteursrecht en licenties in digitale contexten. & Zeer geavanceerd & V
& AI-Impliciet \\
\end{longtable}

\subsubsection{3.4 Computational thinking en
programmeren}\label{computational-thinking-en-programmeren}

\begin{longtable}[]{@{}
  >{\centering\arraybackslash}p{(\linewidth - 8\tabcolsep) * \real{0.1053}}
  >{\centering\arraybackslash}p{(\linewidth - 8\tabcolsep) * \real{0.3947}}
  >{\centering\arraybackslash}p{(\linewidth - 8\tabcolsep) * \real{0.2105}}
  >{\centering\arraybackslash}p{(\linewidth - 8\tabcolsep) * \real{0.1842}}
  >{\centering\arraybackslash}p{(\linewidth - 8\tabcolsep) * \real{0.1053}}@{}}
\toprule\noalign{}
\begin{minipage}[b]{\linewidth}\centering
ID
\end{minipage} & \begin{minipage}[b]{\linewidth}\centering
Leerresultaat
\end{minipage} & \begin{minipage}[b]{\linewidth}\centering
Niveau
\end{minipage} & \begin{minipage}[b]{\linewidth}\centering
K/V/H
\end{minipage} & \begin{minipage}[b]{\linewidth}\centering
AI
\end{minipage} \\
\midrule\noalign{}
\endhead
\bottomrule\noalign{}
\endlastfoot
LO3.4.01 & De essentiële rol van mensen inzien bij het bepalen hoe
computerprogramma's en AI-systemen worden gebruikt. & Basis & H &
AI-Expliciet \\
LO3.4.02 & Veelvoorkomende toepassingen van computerprogramma's en
applicaties identificeren. & Basis & K & AI-Impliciet \\
LO3.4.03 & Computational thinking herkennen als een menselijke
activiteit die het identificeren van stappen omvat die door een computer
kunnen worden uitgevoerd om een probleem of taak op te lossen. & Basis &
K & - \\
LO3.4.04 & In algemene termen herkennen wat AI is. & Basis & K &
AI-Expliciet \\
LO3.4.05 & Op een algemene manier identificeren wat wel en wat niet een
AI-systeem is. & Basis & K & AI-Expliciet \\
LO3.4.06 & Veelvoorkomende voorbeelden van toepassingen van AI-systemen
identificeren. & Basis & K & AI-Expliciet \\
LO3.4.07 & Basisinstructies aan een computer geven om eenvoudige taken
uit te voeren. & Basis & V & AI-Impliciet \\
LO3.4.08 & Eenvoudige reeksen symbolisch weergeven en eenvoudige
symbolische reeksen interpreteren. & Basis & V & - \\
LO3.4.09 & De relevantie inzien van computational thinking,
algoritmische weergave en programmeren in alledaagse contexten. &
Gemiddeld & H & AI-Impliciet \\
LO3.4.10 & Het belang inzien van ethiek en toegankelijkheid in
programmeercontexten. & Gemiddeld & H & AI-Impliciet \\
LO3.4.11 & Onderscheid maken tussen een computationeel model van de
werkelijkheid en de werkelijkheid zelf. & Gemiddeld & K &
AI-Impliciet \\
LO3.4.12 & Het concept van een algoritme herkennen, met voorbeelden uit
computational thinking of programmeren. & Gemiddeld & K & - \\
LO3.4.13 & Verschillen definiëren tussen een berekenbaar probleem en een
niet-berekenbaar probleem. & Gemiddeld & K & - \\
LO3.4.14 & Algemene stappen in computational thinking definiëren. &
Gemiddeld & K & - \\
LO3.4.15 & Herkennen dat er een verscheidenheid aan programmeertalen is,
elk met een reeks mogelijke toepassingen. & Gemiddeld & K & - \\
LO3.4.16 & Fundamentele programmeerconcepten en algemene stappen bij het
programmeren definiëren. & Gemiddeld & K & - \\
LO3.4.17 & De rol van programmeren in robotica herkennen. & Gemiddeld &
K & - \\
LO3.4.18 & Herkennen dat machine learning een tak van AI is die
algoritmen in staat stelt om van data te leren en voorspellingen te
doen. & Gemiddeld & K & AI-Expliciet \\
LO3.4.19 & Herkennen dat er stappen gevolgd moeten worden om een
computerprogramma of een AI-systeem te ontwikkelen, te valideren en te
implementeren. & Gemiddeld & K & AI-Expliciet \\
LO3.4.20 & Voorbeelden van toepassingen van machine learning
beschrijven. & Gemiddeld & K & AI-Expliciet \\
LO3.4.21 & Voorbeelden van AI-systeemtoepassingen uit verschillende
sectoren van de samenleving beschrijven. & Gemiddeld & K &
AI-Expliciet \\
LO3.4.22 & Basisinformatie vertalen naar logische bewerkingen. &
Gemiddeld & V & AI-Impliciet \\
LO3.4.23 & Eenvoudige programma's met controlestructuren ontwikkelen. &
Gemiddeld & V & AI-Impliciet \\
LO3.4.24 & Visuele representaties zoals stroomdiagrammen creëren om
basisalgoritmen te illustreren. & Gemiddeld & V & AI-Impliciet \\
LO3.4.25 & Het belang inzien van menselijk toezicht en mensgerichte
benaderingen bij de ontwikkeling en implementatie van
computerprogramma's en AI-systemen. & Geavanceerd & H & AI-Expliciet \\
LO3.4.26 & De concepten en rol van mensgerichte benaderingen en
menselijk toezicht definiëren in de context van programmeren en
AI-systemen. & Geavanceerd & K & AI-Expliciet \\
LO3.4.27 & De belangrijkste stappen beschrijven bij het ontwikkelen,
valideren en implementeren van een computerprogramma of een AI-systeem.
& Geavanceerd & K & AI-Expliciet \\
LO3.4.28 & Onderscheid maken tussen de belangrijkste soorten machine
learning. & Geavanceerd & K & AI-Expliciet \\
LO3.4.29 & De belangrijkste kenmerken en doelen van veelgebruikte
machine learning-algoritmen identificeren. & Geavanceerd & K &
AI-Expliciet \\
LO3.4.30 & De rol van gebruikerservaring (UX) en klantervaring (CX) bij
programmeren beschrijven. & Geavanceerd & K & - \\
LO3.4.31 & Voorbeelden beschrijven van de toepassing van computational
thinking en programmeren in de robotica. & Geavanceerd & K &
AI-Impliciet \\
LO3.4.32 & Routinetaken identificeren die (gedeeltelijk of volledig)
geautomatiseerd kunnen worden via programmeertools of AI-systemen. &
Geavanceerd & K & AI-Expliciet \\
LO3.4.33 & Ethische en praktische aspecten van de ontwikkeling en
implementatie van computerprogramma's en AI-systemen beoordelen. &
Geavanceerd & V & AI-Expliciet \\
LO3.4.34 & Computational thinking, kennis van programmeren en/of
AI-systemen toepassen om routinetaken (gedeeltelijk of volledig) te
automatiseren. & Geavanceerd & V & AI-Expliciet \\
LO3.4.35 & Programmeerhulpmiddelen of AI-systemen toepassen op complexe
taken op het gebied van computational thinking. & Geavanceerd & V &
AI-Expliciet \\
LO3.4.36 & Ethische programmeer- en/of AI-systeemontwikkelingspraktijken
bevorderen en ondersteunen. & Zeer geavanceerd & H & AI-Expliciet \\
LO3.4.37 & Op de hoogte blijven van de huidige ontwikkelingen in
programmeertechnieken en gerelateerde toepassingen van AI-systemen,
zoals robotica. & Zeer geavanceerd & H & AI-Expliciet \\
LO3.4.38 & Leiding geven aan of bijdragen aan complexe projecten gericht
op toepassingen van computational thinking, programmeren of AI-systemen,
inclusief het ontwikkelen, valideren en implementeren van
computerprogramma's of AI-systemen. & Zeer geavanceerd & V &
AI-Expliciet \\
LO3.4.39 & Anderen helpen bij het ontwikkelen van
basisprogrammeervaardigheden en/of vaardigheden in het toepassen van
AI-systemen bij taken op het gebied van computational thinking. & Zeer
geavanceerd & V & AI-Expliciet \\
\end{longtable}

\subsubsection{4.1 Apparaten beschermen}\label{apparaten-beschermen}

\begin{longtable}[]{@{}
  >{\centering\arraybackslash}p{(\linewidth - 8\tabcolsep) * \real{0.1053}}
  >{\centering\arraybackslash}p{(\linewidth - 8\tabcolsep) * \real{0.3947}}
  >{\centering\arraybackslash}p{(\linewidth - 8\tabcolsep) * \real{0.2105}}
  >{\centering\arraybackslash}p{(\linewidth - 8\tabcolsep) * \real{0.1842}}
  >{\centering\arraybackslash}p{(\linewidth - 8\tabcolsep) * \real{0.1053}}@{}}
\toprule\noalign{}
\begin{minipage}[b]{\linewidth}\centering
ID
\end{minipage} & \begin{minipage}[b]{\linewidth}\centering
Leerresultaat
\end{minipage} & \begin{minipage}[b]{\linewidth}\centering
Niveau
\end{minipage} & \begin{minipage}[b]{\linewidth}\centering
K/V/H
\end{minipage} & \begin{minipage}[b]{\linewidth}\centering
AI
\end{minipage} \\
\midrule\noalign{}
\endhead
\bottomrule\noalign{}
\endlastfoot
LO4.1.01 & Het belang inzien van de eigen individuele rol bij het
beschermen van digitale apparaten en hun inhoud. & Basis & H & - \\
LO4.1.02 & De concepten van cybersecurity, cyberdreiging en cyberaanval
herkennen. & Basis & K & - \\
LO4.1.03 & Herkennen dat de acties van individuen en
cybersecurityhulpmiddelen samenwerken om apparaten en hun inhoud veilig
te houden. & Basis & K & - \\
LO4.1.04 & Herkennen dat er cybersecuritywetgeving bestaat die helpt de
veiligheid van producten en diensten te waarborgen. & Basis & K & - \\
LO4.1.05 & Basisbeschermingsmaatregelen en -praktijken voor apparaten
identificeren, zoals antivirussoftware, schermvergrendeling, sterke
wachtwoorden en multifactorauthenticatie. & Basis & K & - \\
LO4.1.06 & Basisbeschermingsmaatregelen en -praktijken voor apparaten
toepassen, zoals antivirussoftware, schermvergrendeling, sterke
wachtwoorden en multifactorauthenticatie. & Basis & V & - \\
LO4.1.07 & Het belang inzien van waakzaam blijven en op de hoogte
blijven van cybersecuritypraktijken. & Gemiddeld & H & AI-Impliciet \\
LO4.1.08 & Belangrijkste kenmerken van malware beschrijven. & Gemiddeld
& K & AI-Impliciet \\
LO4.1.09 & Herkennen dat AI-systemen kunnen worden gebruikt voor zowel
cyberaanvallen als cybersecurity. & Gemiddeld & K & AI-Expliciet \\
LO4.1.10 & Verschillende technieken en praktijken voor malwarepreventie
toepassen om apparaten en hun inhoud te beschermen. & Gemiddeld & V &
AI-Impliciet \\
LO4.1.11 & Prioriteit geven aan het regelmatig controleren en bijwerken
van cybersecuritymaatregelen op de eigen apparaten. & Geavanceerd & H &
- \\
LO4.1.12 & Belangrijke rechten van individuen onder de huidige
cybersecuritywetgeving beschrijven. & Geavanceerd & K & AI-Impliciet \\
LO4.1.13 & Voorbeelden identificeren van hoe recente en opkomende
technologieën zoals AI-systemen worden gebruikt bij cyberaanvallen en
cybersecurity. & Geavanceerd & K & AI-Expliciet \\
LO4.1.14 & Cybersecuritymaatregelen bijwerken om apparaten en hun inhoud
te beschermen als reactie op evoluerende digitale dreigingen. &
Geavanceerd & V & AI-Impliciet \\
LO4.1.15 & Anderen helpen bij het implementeren van
basisbeschermingsmaatregelen voor apparaten, zoals antivirussoftware,
schermvergrendeling, sterke wachtwoorden en multifactorauthenticatie. &
Geavanceerd & V & - \\
LO4.1.16 & Op de hoogte blijven van technologische en wettelijke
ontwikkelingen op het gebied van cybersecurity. & Zeer geavanceerd & H &
AI-Impliciet \\
LO4.1.17 & Rechten van individuen beoordelen onder relevante bepalingen
van de huidige cybersecuritywetgeving. & Zeer geavanceerd & V &
AI-Impliciet \\
LO4.1.18 & Leiding geven aan of bijdragen aan op de burger gerichte
cybersecurity-initiatieven. & Zeer geavanceerd & V & AI-Impliciet \\
LO4.1.19 & Anderen ondersteunen bij het opbouwen van hun vaardigheden om
hun apparaten en hun inhoud te beschermen tegen digitale dreigingen. &
Zeer geavanceerd & V & AI-Impliciet \\
\end{longtable}

\subsubsection{4.2 Persoonsgegevens en privacy
beschermen}\label{persoonsgegevens-en-privacy-beschermen}

\begin{longtable}[]{@{}
  >{\centering\arraybackslash}p{(\linewidth - 8\tabcolsep) * \real{0.1053}}
  >{\centering\arraybackslash}p{(\linewidth - 8\tabcolsep) * \real{0.3947}}
  >{\centering\arraybackslash}p{(\linewidth - 8\tabcolsep) * \real{0.2105}}
  >{\centering\arraybackslash}p{(\linewidth - 8\tabcolsep) * \real{0.1842}}
  >{\centering\arraybackslash}p{(\linewidth - 8\tabcolsep) * \real{0.1053}}@{}}
\toprule\noalign{}
\begin{minipage}[b]{\linewidth}\centering
ID
\end{minipage} & \begin{minipage}[b]{\linewidth}\centering
Leerresultaat
\end{minipage} & \begin{minipage}[b]{\linewidth}\centering
Niveau
\end{minipage} & \begin{minipage}[b]{\linewidth}\centering
K/V/H
\end{minipage} & \begin{minipage}[b]{\linewidth}\centering
AI
\end{minipage} \\
\midrule\noalign{}
\endhead
\bottomrule\noalign{}
\endlastfoot
LO4.2.01 & Het belang inzien van een voorzichtige aanpak bij het delen
van persoonlijke data in digitale omgevingen. & Basis & H & - \\
LO4.2.02 & Herkennen dat persoonlijke data wordt verzameld en
gegenereerd via een grote verscheidenheid aan digitale bronnen en
processen. & Basis & K & AI-Impliciet \\
LO4.2.03 & Risico's identificeren van het delen van persoonlijke data in
digitale omgevingen, inclusief specifieke risico's in relatie tot
AI-systemen. & Basis & K & AI-Expliciet \\
LO4.2.04 & Herkennen dat individuen recht hebben op privacy en dat hun
persoonlijke data beschermd is onder de wetgeving. & Basis & K &
AI-Impliciet \\
LO4.2.05 & Herkennen dat manipulatieve methoden kunnen worden gebruikt
om individuen te misleiden om toegang te verlenen tot persoonlijke data,
accounts of andere gevoelige informatie. & Basis & K & - \\
LO4.2.06 & Tekenen van identiteitsfraude herkennen. & Basis & K & - \\
LO4.2.07 & Herkennen dat gebruikers van platforms en diensten kunnen
verzoeken om informatie die ongepast online is gedeeld te blokkeren of
te markeren. & Basis & K & - \\
LO4.2.08 & Basisbeveiligingsmaatregelen implementeren voor online
betalingen en transacties. & Basis & V & - \\
LO4.2.09 & Passend reageren op tekenen van identiteitsfraude. & Basis &
V & - \\
LO4.2.10 & Persoonlijke informatie blokkeren of markeren die op
ongepaste wijze online is gedeeld. & Basis & V & - \\
LO4.2.11 & Het belang inzien van zorgvuldige omgang met persoonlijke
data van zichzelf en anderen, vooral van kwetsbare individuen en
kinderen. & Gemiddeld & H & - \\
LO4.2.12 & Belangrijke concepten van gegevensbescherming en
privacywetgeving herkennen, waaronder privacy, anonimisering,
pseudonimisering en het recht op gegevensverwijdering. & Gemiddeld & K &
AI-Impliciet \\
LO4.2.13 & Het doel van privacyverklaringen identificeren. & Gemiddeld &
K & - \\
LO4.2.14 & Belangrijke concepten uit privacyverklaringen definiëren,
zoals betrokkene, bewaartermijn, gegevensoverdracht en geautomatiseerd
besluitvormingssysteem. & Gemiddeld & K & - \\
LO4.2.15 & Technieken gerelateerd aan social engineering in digitale
omgevingen beschrijven, zoals phishing of baiting. & Gemiddeld & K &
- \\
LO4.2.16 & Een inbreuk in verband met persoonsgegevens (datalek)
definiëren onder de huidige gegevensbescherming- en privacywetgeving. &
Gemiddeld & K & - \\
LO4.2.17 & Herkennen dat de regulering van het eigendom van persoonlijke
data in AI-systemen complex is. & Gemiddeld & K & AI-Expliciet \\
LO4.2.18 & Privacy-implicaties beschrijven die verbonden zijn aan het
gebruik van online gedeelde inhoud, bijvoorbeeld om AI-systemen te
trainen. & Gemiddeld & K & AI-Expliciet \\
LO4.2.19 & Belangrijkste kenmerken en functies van privacyhulpmiddelen
definiëren. & Gemiddeld & K & AI-Impliciet \\
LO4.2.20 & Persoonlijke data en privacy beheren in verschillende
digitale omgevingen, inclusief het gebruik van privacyhulpmiddelen. &
Gemiddeld & V & AI-Impliciet \\
LO4.2.21 & Voortdurend kwesties rondom data-eigendom en privacy
verkennen in relatie tot digitale technologische ontwikkelingen. &
Geavanceerd & H & AI-Impliciet \\
LO4.2.22 & Anderen ondersteunen bij het begrijpen van hun rechten onder
de huidige gegevensbescherming- en privacywetgeving. & Geavanceerd & V &
AI-Impliciet \\
LO4.2.23 & Anderen helpen bij het implementeren van basisstrategieën om
persoonlijke data te beschermen en privacy in digitale omgevingen te
beheren. & Geavanceerd & V & AI-Impliciet \\
LO4.2.24 & Op de hoogte blijven van technologische en wettelijke
ontwikkelingen op het gebied van persoonlijke data, data-eigendom en
privacy. & Zeer geavanceerd & H & AI-Impliciet \\
LO4.2.25 & Adviseren over beleidsmatige of regelgevende aspecten van
gegevensbescherming en privacy in digitale contexten. & Zeer geavanceerd
& V & AI-Impliciet \\
LO4.2.26 & Leiding geven aan of bijdragen aan het ontwerp van
strategieën voor de bescherming van persoonlijke data in digitale
contexten. & Zeer geavanceerd & V & AI-Impliciet \\
\end{longtable}

\subsubsection{4.3 Welzijn ondersteunen}\label{welzijn-ondersteunen}

\begin{longtable}[]{@{}
  >{\centering\arraybackslash}p{(\linewidth - 8\tabcolsep) * \real{0.1053}}
  >{\centering\arraybackslash}p{(\linewidth - 8\tabcolsep) * \real{0.3947}}
  >{\centering\arraybackslash}p{(\linewidth - 8\tabcolsep) * \real{0.2105}}
  >{\centering\arraybackslash}p{(\linewidth - 8\tabcolsep) * \real{0.1842}}
  >{\centering\arraybackslash}p{(\linewidth - 8\tabcolsep) * \real{0.1053}}@{}}
\toprule\noalign{}
\begin{minipage}[b]{\linewidth}\centering
ID
\end{minipage} & \begin{minipage}[b]{\linewidth}\centering
Leerresultaat
\end{minipage} & \begin{minipage}[b]{\linewidth}\centering
Niveau
\end{minipage} & \begin{minipage}[b]{\linewidth}\centering
K/V/H
\end{minipage} & \begin{minipage}[b]{\linewidth}\centering
AI
\end{minipage} \\
\midrule\noalign{}
\endhead
\bottomrule\noalign{}
\endlastfoot
LO4.3.01 & De voordelen inzien van een balans tussen online en offline
activiteiten. & Basis & H & - \\
LO4.3.02 & Het belang inzien van dagelijkse routines die digitale stress
minimaliseren en sociale verbondenheid bevorderen. & Basis & H & - \\
LO4.3.03 & Belangrijkste risico's en voordelen voor het fysieke, mentale
en sociale welzijn in digitale omgevingen identificeren. & Basis & K &
AI-Impliciet \\
LO4.3.04 & Herkennen dat er een verscheidenheid aan informatie, groepen
en gemeenschappen in digitale omgevingen is die het eigen fysieke,
mentale en/of sociale welzijn kunnen ondersteunen. & Basis & K &
AI-Impliciet \\
LO4.3.05 & Risico's en voordelen herkennen voor het eigen fysieke,
mentale en sociale welzijn in digitale omgevingen. & Basis & K &
AI-Impliciet \\
LO4.3.06 & Kenmerken identificeren van digitale platforms of diensten,
zoals sociale media, die ontworpen zijn om de aandacht van individuen te
trekken en vast te houden. & Basis & K & AI-Impliciet \\
LO4.3.07 & Beperkingen en risico's identificeren van het gebruik van
virtuele assistenten en AI-systemen om menselijk welzijn te
ondersteunen. & Basis & K & AI-Expliciet \\
LO4.3.08 & Strategieën identificeren om fysiek, mentaal en sociaal
welzijn in digitale omgevingen te ondersteunen. & Basis & K & - \\
LO4.3.09 & Tekenen en mogelijke effecten van problematisch gebruik van
digitale technologieën herkennen. & Basis & K & - \\
LO4.3.10 & Herkennen dat er wetten en regels zijn die helpen het welzijn
van individuen in digitale omgevingen te beschermen. & Basis & K &
AI-Impliciet \\
LO4.3.11 & Een basisbeoordeling maken van de eigen digitale gewoonten in
relatie tot het eigen fysieke, mentale en sociale welzijn. & Basis & V &
- \\
LO4.3.12 & Gepersonaliseerde strategieën toepassen om fysiek, mentaal en
sociaal welzijn in digitale omgevingen te ondersteunen. & Basis & V &
- \\
LO4.3.13 & Het belang inzien van het recht van zichzelf en anderen om
offline te zijn (onbereikbaar te zijn). & Gemiddeld & H & - \\
LO4.3.14 & De fysieke, mentale en sociale voordelen inzien van het
analyseren van de eigen digitale gebruikspatronen. & Gemiddeld & H &
- \\
LO4.3.15 & Betrouwbare informatiebronnen en inclusieve groepen en
gemeenschappen in digitale omgevingen identificeren die het eigen
fysieke, mentale en/of sociale welzijn kunnen ondersteunen. & Gemiddeld
& K & AI-Impliciet \\
LO4.3.16 & Voorbeelden beschrijven van schadelijke inhoud en gedrag in
digitale omgevingen en de mogelijke effecten daarvan op zichzelf en
anderen. & Gemiddeld & K & AI-Impliciet \\
LO4.3.17 & Manieren beschrijven waarop sommige toepassingen van digitale
technologieën, zoals sociale media, bias, stereotypering en uitsluiting
versterken en in stand houden. & Gemiddeld & K & AI-Impliciet \\
LO4.3.18 & Strategieën beschrijven om te helpen beschermen tegen en
effectief te reageren op schadelijk gedrag en inhoud. & Gemiddeld & K &
AI-Impliciet \\
LO4.3.19 & Mogelijke manieren identificeren om schadelijk gedrag of
inhoud te melden of in te grijpen als dit in digitale omgevingen wordt
aangetroffen. & Gemiddeld & K & AI-Impliciet \\
LO4.3.20 & De effecten van schadelijk gedrag, inhoud en misleidend
ontwerp (deceptive design) in digitale omgevingen op zichzelf en anderen
beschrijven. & Gemiddeld & K & AI-Impliciet \\
LO4.3.21 & De eigen digitale gebruikspatronen analyseren en aanpassen
ter ondersteuning van het fysieke, mentale en sociale welzijn. &
Gemiddeld & V & - \\
LO4.3.22 & Strategieën implementeren om zichzelf te helpen beschermen
tegen en effectief te reageren op schadelijk gedrag en inhoud. &
Gemiddeld & V & - \\
LO4.3.23 & Aanpassen aan veranderende digitale technologische
ontwikkelingen en behoeften om het fysieke, mentale en sociale welzijn
van zichzelf en anderen te ondersteunen en te behouden. & Geavanceerd &
H & AI-Impliciet \\
LO4.3.24 & Voortdurend de rol van digitale technologieën zoals sociale
media onderzoeken bij het versterken en in stand houden van bias,
stereotypering en uitsluiting. & Geavanceerd & H & AI-Impliciet \\
LO4.3.25 & Anderen helpen hun gebruik van digitale technologieën te
herzien en aan te passen ter ondersteuning en behoud van hun fysieke,
mentale en sociale welzijn. & Geavanceerd & V & - \\
LO4.3.26 & Effectief melden van of ingrijpen bij gevallen van schadelijk
gedrag of inhoud in digitale omgevingen. & Geavanceerd & V &
AI-Impliciet \\
LO4.3.27 & Anderen helpen vaardigheden op te bouwen om de rol van
digitale technologieën zoals sociale media bij het versterken en in
stand houden van bias, stereotypering en uitsluiting tegen te gaan. &
Geavanceerd & V & AI-Impliciet \\
LO4.3.28 & Anderen helpen hun rechten te begrijpen met betrekking tot
welzijn en/of inclusie in digitale omgevingen. & Geavanceerd & V &
AI-Impliciet \\
LO4.3.29 & Anderen helpen bewustzijn te ontwikkelen over schadelijk
gedrag, inhoud en misleidend ontwerp in digitale omgevingen. &
Geavanceerd & V & AI-Impliciet \\
LO4.3.30 & Acties bevorderen en ondersteunen die welzijn en inclusie in
digitale omgevingen bevorderen. & Zeer geavanceerd & H & AI-Impliciet \\
LO4.3.31 & Bewijs over welzijn en/of inclusie in digitale omgevingen
beoordelen en evalueren om besluitvorming te sturen. & Zeer geavanceerd
& V & AI-Impliciet \\
LO4.3.32 & Leiding geven aan of bijdragen aan initiatieven die welzijn
en/of inclusie in digitale omgevingen ondersteunen. & Zeer geavanceerd &
V & AI-Impliciet \\
LO4.3.33 & Bijdragen aan juridische en regelgevende besluitvorming met
betrekking tot het welzijn en/of de inclusie van individuen in digitale
omgevingen. & Zeer geavanceerd & V & AI-Impliciet \\
\end{longtable}

\subsubsection{4.4 Milieueffecten van digitale
technologieën}\label{milieueffecten-van-digitale-technologieuxebn}

\begin{longtable}[]{@{}
  >{\centering\arraybackslash}p{(\linewidth - 8\tabcolsep) * \real{0.1053}}
  >{\centering\arraybackslash}p{(\linewidth - 8\tabcolsep) * \real{0.3947}}
  >{\centering\arraybackslash}p{(\linewidth - 8\tabcolsep) * \real{0.2105}}
  >{\centering\arraybackslash}p{(\linewidth - 8\tabcolsep) * \real{0.1842}}
  >{\centering\arraybackslash}p{(\linewidth - 8\tabcolsep) * \real{0.1053}}@{}}
\toprule\noalign{}
\begin{minipage}[b]{\linewidth}\centering
ID
\end{minipage} & \begin{minipage}[b]{\linewidth}\centering
Leerresultaat
\end{minipage} & \begin{minipage}[b]{\linewidth}\centering
Niveau
\end{minipage} & \begin{minipage}[b]{\linewidth}\centering
K/V/H
\end{minipage} & \begin{minipage}[b]{\linewidth}\centering
AI
\end{minipage} \\
\midrule\noalign{}
\endhead
\bottomrule\noalign{}
\endlastfoot
LO4.4.01 & De rol inzien die individuen kunnen spelen om de
milieu-impact van digitale technologieën te helpen verminderen. & Basis
& H & - \\
LO4.4.02 & Herkennen dat sommige digitale technologieën en
infrastructuren, zoals AI-systemen en datacenters, grote effecten hebben
op het milieu. & Basis & K & AI-Expliciet \\
LO4.4.03 & Herkennen dat de volledige milieu-impact van digitale
technologieën niet onmiddellijk zichtbaar is voor een individuele
gebruiker. & Basis & K & AI-Impliciet \\
LO4.4.04 & Herkennen dat digitale technologieën energie-efficiëntie en
duurzaamheid kunnen ondersteunen. & Basis & K & AI-Impliciet \\
LO4.4.05 & Eenvoudige strategieën identificeren om het energie- en
datagebruik van digitale technologieën te verminderen, zoals het
minimaliseren van het gebruik van energie-intensieve applicaties en het
niet gebruiken van digitale technologieën wanneer deze niet nodig zijn.
& Basis & K & AI-Impliciet \\
LO4.4.06 & Eenvoudige strategieën toepassen om het energie- en
datagebruik van digitale technologieën te verminderen, zoals het
minimaliseren van het gebruik van energie-intensieve applicaties en het
niet gebruiken van digitale technologieën wanneer deze niet nodig zijn.
& Basis & V & AI-Impliciet \\
LO4.4.07 & Voortdurend de milieu-impact van het eigen gebruik van
digitale technologieën beoordelen. & Gemiddeld & H & AI-Impliciet \\
LO4.4.08 & Milieu-impact van digitale technologieën identificeren die
optreedt tijdens productie, gebruik en afvalverwerking. & Gemiddeld & K
& AI-Impliciet \\
LO4.4.09 & De milieu-impact van datacenters en e-commerce beschrijven. &
Gemiddeld & K & AI-Impliciet \\
LO4.4.10 & Voorbeelden beschrijven van hoe digitale hulpmiddelen een
duurzame levensstijl kunnen ondersteunen. & Gemiddeld & K &
AI-Impliciet \\
LO4.4.11 & De concepten van de deeleconomie en de circulaire economie
definiëren, inclusief risico's, beperkingen en potentiële
milieuvoordelen. & Gemiddeld & K & - \\
LO4.4.12 & Verschillende strategieën toepassen om de milieu-impact van
het eigen gebruik van digitale technologieën en digitale apparaten te
verminderen, zoals geïnformeerde beslissingen bij de aankoop van
digitale apparaten, recycling en reparatie van apparaten, milieubewuste
e-commercepraktijken en milieubewuste gebruikspatronen. & Gemiddeld & V
& AI-Impliciet \\
LO4.4.13 & Op de hoogte blijven van de milieu-impact van digitale
technologieën en de manieren waarop digitale technologieën duurzaamheid
kunnen ondersteunen. & Geavanceerd & H & AI-Impliciet \\
LO4.4.14 & De milieu-impact van digitale technologieën en
infrastructuren evalueren ter ondersteuning van besluitvorming of
belangenbehartiging. & Geavanceerd & V & AI-Impliciet \\
LO4.4.15 & Anderen helpen hun gebruik van digitale technologieën te
beoordelen om manieren te identificeren waarop de milieu-impact kan
worden verminderd. & Geavanceerd & V & AI-Impliciet \\
LO4.4.16 & Op de hoogte blijven van de implicaties voor milieu en
duurzaamheid van digitale technologieën in verschillende sectoren. &
Zeer geavanceerd & H & AI-Impliciet \\
LO4.4.17 & Acties voor milieuvriendelijk duurzaam gebruik van digitale
technologieën bevorderen en ondersteunen. & Zeer geavanceerd & H &
AI-Impliciet \\
LO4.4.18 & Leiding geven aan of bijdragen aan initiatieven voor digitale
duurzaamheid. & Zeer geavanceerd & V & AI-Impliciet \\
LO4.4.19 & Bijdragen aan verbeteringen in of oplossingen voor digitale
duurzaamheid. & Zeer geavanceerd & V & AI-Impliciet \\
\end{longtable}

\subsubsection{5.1 Technische problemen identificeren en
oplossen}\label{technische-problemen-identificeren-en-oplossen}

\begin{longtable}[]{@{}
  >{\centering\arraybackslash}p{(\linewidth - 8\tabcolsep) * \real{0.1053}}
  >{\centering\arraybackslash}p{(\linewidth - 8\tabcolsep) * \real{0.3947}}
  >{\centering\arraybackslash}p{(\linewidth - 8\tabcolsep) * \real{0.2105}}
  >{\centering\arraybackslash}p{(\linewidth - 8\tabcolsep) * \real{0.1842}}
  >{\centering\arraybackslash}p{(\linewidth - 8\tabcolsep) * \real{0.1053}}@{}}
\toprule\noalign{}
\begin{minipage}[b]{\linewidth}\centering
ID
\end{minipage} & \begin{minipage}[b]{\linewidth}\centering
Leerresultaat
\end{minipage} & \begin{minipage}[b]{\linewidth}\centering
Niveau
\end{minipage} & \begin{minipage}[b]{\linewidth}\centering
K/V/H
\end{minipage} & \begin{minipage}[b]{\linewidth}\centering
AI
\end{minipage} \\
\midrule\noalign{}
\endhead
\bottomrule\noalign{}
\endlastfoot
LO5.1.01 & Inzien dat technische problemen in digitale omgevingen veel
voorkomen. & Basis & H & - \\
LO5.1.02 & De voordelen inzien van het zoeken naar hulp bij het oplossen
van technische problemen. & Basis & H & - \\
LO5.1.03 & Onderscheid maken tussen besturingssystemen en software. &
Basis & K & - \\
LO5.1.04 & Belangrijkste kenmerken van hardware, software,
connectiviteit en veelvoorkomende randapparatuur identificeren. & Basis
& K & - \\
LO5.1.05 & Tekenen van veelvoorkomende technische problemen herkennen,
zoals gebrek aan connectiviteit, vergeten wachtwoord, vergeten
bestandslocatie, niet-opgeslagen document of fout in e-mail- of
webadres. & Basis & K & - \\
LO5.1.06 & Instructies volgen om veelvoorkomende technische problemen op
te lossen, zoals gebrek aan connectiviteit, vergeten wachtwoord,
vergeten bestandslocatie, niet-opgeslagen document of fout in e-mail- of
webadres. & Basis & V & - \\
LO5.1.07 & Indien nodig software en applicaties installeren en
bijwerken. & Basis & V & - \\
LO5.1.08 & De voordelen inzien van het opbouwen van vaardigheden en
autonomie bij het aanpakken van technische problemen. & Gemiddeld & H &
- \\
LO5.1.09 & Technische problemen in digitale omgevingen oplossen met
behulp van verschillende zoek- en probleemoplossingsstrategieën (met
menselijke hulp of met hulp van digitale technologie). & Gemiddeld & V &
AI-Impliciet \\
LO5.1.10 & Instellingen op hoofd- en randapparatuur bijwerken en
aanpassen om goede prestaties te behouden. & Gemiddeld & V & - \\
LO5.1.11 & Prioriteit geven aan de ontwikkeling van het eigen vermogen
om technische problemen in digitale omgevingen te diagnosticeren en op
te lossen. & Geavanceerd & H & - \\
LO5.1.12 & Anderen helpen bij het diagnosticeren en oplossen van
technische problemen in digitale omgevingen. & Geavanceerd & V & - \\
LO5.1.13 & Verschillende strategieën voor het vinden van oplossingen
gebruiken voor het verhelpen van complexe technische problemen in
digitale omgevingen. & Geavanceerd & V & AI-Impliciet \\
LO5.1.14 & Anderen helpen bij het ontwikkelen van vertrouwen en
autonomie om technische problemen in digitale omgevingen op te lossen. &
Zeer geavanceerd & H & - \\
LO5.1.15 & Training ontwerpen of geven ter ondersteuning van het gebruik
van digitale apparaten of systemen. & Zeer geavanceerd & V &
AI-Impliciet \\
\end{longtable}

\subsubsection{5.2 Behoeften en digitale technologische antwoorden
identificeren}\label{behoeften-en-digitale-technologische-antwoorden-identificeren}

\begin{longtable}[]{@{}
  >{\centering\arraybackslash}p{(\linewidth - 8\tabcolsep) * \real{0.1053}}
  >{\centering\arraybackslash}p{(\linewidth - 8\tabcolsep) * \real{0.3947}}
  >{\centering\arraybackslash}p{(\linewidth - 8\tabcolsep) * \real{0.2105}}
  >{\centering\arraybackslash}p{(\linewidth - 8\tabcolsep) * \real{0.1842}}
  >{\centering\arraybackslash}p{(\linewidth - 8\tabcolsep) * \real{0.1053}}@{}}
\toprule\noalign{}
\begin{minipage}[b]{\linewidth}\centering
ID
\end{minipage} & \begin{minipage}[b]{\linewidth}\centering
Leerresultaat
\end{minipage} & \begin{minipage}[b]{\linewidth}\centering
Niveau
\end{minipage} & \begin{minipage}[b]{\linewidth}\centering
K/V/H
\end{minipage} & \begin{minipage}[b]{\linewidth}\centering
AI
\end{minipage} \\
\midrule\noalign{}
\endhead
\bottomrule\noalign{}
\endlastfoot
LO5.2.01 & Het belang inzien van individuele keuzes bij de configuratie
van digitale omgevingen. & Basis & H & - \\
LO5.2.02 & Veelvoorkomende manieren herkennen waarop kenmerken van
digitale omgevingen kunnen worden aangepast aan de behoeften en
voorkeuren van gebruikers. & Basis & K & - \\
LO5.2.03 & Het concept en het doel van een digitaal
assistentiehulpmiddel herkennen. & Basis & K & AI-Impliciet \\
LO5.2.04 & De aanwezigheid van AI-systemen in digitale
assistentiehulpmiddelen herkennen. & Basis & K & AI-Expliciet \\
LO5.2.05 & Veelvoorkomende ondersteunende technologieën en hun doelen
identificeren. & Basis & K & AI-Impliciet \\
LO5.2.06 & Ondersteunende technologieën gebruiken indien en voor zover
nodig. & Basis & V & AI-Impliciet \\
LO5.2.07 & Digitale assistentiehulpmiddelen gebruiken om eenvoudige
taken te ondersteunen, met bewustzijn van hun voordelen en beperkingen.
& Basis & V & AI-Impliciet \\
LO5.2.08 & De voordelen inzien van het verkennen van aanpassingen aan
configuraties van digitale omgevingen en kenmerken van digitale
assistentiehulpmiddelen. & Gemiddeld & H & AI-Impliciet \\
LO5.2.09 & Kenmerken van de eigen digitale omgeving aanpassen aan de
behoeften en voorkeuren van zichzelf en anderen. & Gemiddeld & V &
AI-Impliciet \\
LO5.2.10 & Geïnformeerd gebruik maken van digitale
assistentiehulpmiddelen om aan de eigen behoeften en die van anderen te
voldoen, met bewustzijn van hun voordelen en beperkingen. & Gemiddeld &
V & AI-Impliciet \\
LO5.2.11 & Prioriteit geven aan een voortdurende beoordeling van hoe
configuraties van de digitale omgeving, digitale assistentiehulpmiddelen
en/of ondersteunende technologieën kunnen voldoen aan de behoeften van
zichzelf en anderen. & Geavanceerd & H & AI-Impliciet \\
LO5.2.12 & Kenmerken van de eigen digitale omgeving aanpassen aan de
behoeften en voorkeuren van zichzelf en anderen. & Geavanceerd & V &
AI-Impliciet \\
LO5.2.13 & De toegankelijkheid, inclusiviteit, eerlijkheid en/of
rechtgevoeligheid van digitale technologieën in een gegeven context
beoordelen. & Geavanceerd & V & AI-Impliciet \\
LO5.2.14 & Anderen ondersteunen bij het maken van geïnformeerd gebruik
van digitale assistentiehulpmiddelen en aanpassingen aan configuraties
van de digitale omgeving om aan hun eigen behoeften en die van anderen
te voldoen. & Geavanceerd & V & AI-Impliciet \\
LO5.2.15 & Het gebruik van inclusieve en toegankelijke digitale
technologieën bevorderen en ondersteunen. & Zeer geavanceerd & H &
AI-Impliciet \\
LO5.2.16 & Complexe behoeften van individuen beoordelen om op maat
gemaakte digitale oplossingen te identificeren en/of te ontwerpen. &
Zeer geavanceerd & V & AI-Impliciet \\
LO5.2.17 & Bijdragen aan verbeteringen in digitale
assistentiehulpmiddelen, toegankelijke configuraties van de digitale
omgeving en/of ondersteunende technologieën. & Zeer geavanceerd & V &
AI-Impliciet \\
\end{longtable}

\subsubsection{5.3 Creatieve oplossingen identificeren met behulp van
digitale
technologieën}\label{creatieve-oplossingen-identificeren-met-behulp-van-digitale-technologieuxebn}

\begin{longtable}[]{@{}
  >{\centering\arraybackslash}p{(\linewidth - 8\tabcolsep) * \real{0.1053}}
  >{\centering\arraybackslash}p{(\linewidth - 8\tabcolsep) * \real{0.3947}}
  >{\centering\arraybackslash}p{(\linewidth - 8\tabcolsep) * \real{0.2105}}
  >{\centering\arraybackslash}p{(\linewidth - 8\tabcolsep) * \real{0.1842}}
  >{\centering\arraybackslash}p{(\linewidth - 8\tabcolsep) * \real{0.1053}}@{}}
\toprule\noalign{}
\begin{minipage}[b]{\linewidth}\centering
ID
\end{minipage} & \begin{minipage}[b]{\linewidth}\centering
Leerresultaat
\end{minipage} & \begin{minipage}[b]{\linewidth}\centering
Niveau
\end{minipage} & \begin{minipage}[b]{\linewidth}\centering
K/V/H
\end{minipage} & \begin{minipage}[b]{\linewidth}\centering
AI
\end{minipage} \\
\midrule\noalign{}
\endhead
\bottomrule\noalign{}
\endlastfoot
LO5.3.01 & Inzien dat digitale technologieën de menselijke creativiteit
kunnen ondersteunen, maar niet vervangen. & Basis & H & AI-Impliciet \\
LO5.3.02 & Voorbeelden identificeren van hoe digitale technologieën
worden gebruikt om real-world problemen op te lossen en verbeteringen
aan te brengen in of nieuwe oplossingen te creëren voor producten en
diensten. & Basis & K & AI-Impliciet \\
LO5.3.03 & Voorbeelden identificeren van waar digitale technologieën
menselijke creativiteit kunnen ondersteunen of versterken. & Basis & K &
AI-Impliciet \\
LO5.3.04 & Het belang inzien van het centraal stellen van mensenrechten,
waarden, behoeften en ervaringen bij het ontwerp en gebruik van digitale
technologieën. & Gemiddeld & H & AI-Impliciet \\
LO5.3.05 & Het concept van mensgerichtheid definiëren en de rol ervan
bij de ontwikkeling en het gebruik van digitale technologieën. &
Gemiddeld & K & AI-Impliciet \\
LO5.3.06 & Voorbeelden identificeren van de wisselwerking tussen mensen
en digitale technologieën bij creativiteit en probleemoplossing. &
Gemiddeld & K & AI-Impliciet \\
LO5.3.07 & Sterktes, zwaktes en ethische overwegingen van digitale
technologieën, inclusief AI-systemen, beschrijven in relatie tot
menselijke creativiteit en probleemoplossing. & Gemiddeld & K &
AI-Expliciet \\
LO5.3.08 & Analyseer de sterktes en zwaktes van beschikbare digitale
technologieën met betrekking tot een specifieke probleemoplossende taak.
& Gemiddeld & V & AI-Impliciet \\
LO5.3.09 & Op een verantwoorde en ethische manier gebruikmaken van
verschillende digitale technologieën om probleemoplossing als individu
of in een groep te ondersteunen. & Gemiddeld & V & AI-Impliciet \\
LO5.3.10 & Prioriteit geven aan mensgerichte benaderingen bij het eigen
gebruik van digitale technologieën voor probleemoplossing. & Geavanceerd
& H & AI-Impliciet \\
LO5.3.11 & Verschillende digitale technologieën efficiënt, verantwoord
en ethisch gebruiken om complexe problemen te helpen oplossen. &
Geavanceerd & V & AI-Impliciet \\
LO5.3.12 & Bijdragen aan initiatieven gericht op de toepassing van
digitale technologieën om complexe probleemoplossende taken te helpen
uitvoeren. & Geavanceerd & V & AI-Impliciet \\
LO5.3.13 & Anderen ondersteunen bij het ontwikkelen van hun vertrouwen
en vaardigheden in het gebruik van digitale technologieën om
praktijkproblemen te helpen oplossen. & Geavanceerd & V &
AI-Impliciet \\
LO5.3.14 & Leiding geven aan of bijdragen aan initiatieven gericht op de
toepassing van digitale technologieën voor zeer complexe of
gespecialiseerde probleemoplossing. & Zeer geavanceerd & V &
AI-Impliciet \\
LO5.3.15 & Leiding geven aan of bijdragen aan initiatieven die digitale
technologieën gebruiken om verbeteringen aan te brengen in of nieuwe
oplossingen te vinden voor praktijkproblemen. & Zeer geavanceerd & V &
AI-Impliciet \\
LO5.3.16 & Anderen ondersteunen bij het ontwikkelen van hun vaardigheden
om digitale technologieën te gebruiken voor complexe of gespecialiseerde
probleemoplossende taken. & Zeer geavanceerd & V & AI-Impliciet \\
\end{longtable}

\subsubsection{5.4 Behoeften aan digitale competenties identificeren en
aanpakken}\label{behoeften-aan-digitale-competenties-identificeren-en-aanpakken}

\begin{longtable}[]{@{}
  >{\centering\arraybackslash}p{(\linewidth - 8\tabcolsep) * \real{0.1053}}
  >{\centering\arraybackslash}p{(\linewidth - 8\tabcolsep) * \real{0.3947}}
  >{\centering\arraybackslash}p{(\linewidth - 8\tabcolsep) * \real{0.2105}}
  >{\centering\arraybackslash}p{(\linewidth - 8\tabcolsep) * \real{0.1842}}
  >{\centering\arraybackslash}p{(\linewidth - 8\tabcolsep) * \real{0.1053}}@{}}
\toprule\noalign{}
\begin{minipage}[b]{\linewidth}\centering
ID
\end{minipage} & \begin{minipage}[b]{\linewidth}\centering
Leerresultaat
\end{minipage} & \begin{minipage}[b]{\linewidth}\centering
Niveau
\end{minipage} & \begin{minipage}[b]{\linewidth}\centering
K/V/H
\end{minipage} & \begin{minipage}[b]{\linewidth}\centering
AI
\end{minipage} \\
\midrule\noalign{}
\endhead
\bottomrule\noalign{}
\endlastfoot
LO5.4.01 & De waarde inzien van het leren over digitale technologieën in
relatie tot de eigen interesses en behoeften. & Basis & H &
AI-Impliciet \\
LO5.4.02 & De voordelen en de normaliteit inzien van het zoeken naar
ondersteuning bij het aanpakken van behoeften op het gebied van digitale
competenties. & Basis & H & - \\
LO5.4.03 & Herkennen dat digitale competentie veel breder is dan alleen
technische vaardigheden. & Basis & K & AI-Impliciet \\
LO5.4.04 & Herkennen dat digitale competentie regelmatige bijscholing
vereist voor het dagelijks leven, werk en leren. & Basis & K &
AI-Impliciet \\
LO5.4.05 & Mogelijkheden identificeren om de eigen digitale competenties
te verbeteren. & Basis & K & AI-Impliciet \\
LO5.4.06 & De voordelen inzien van het op de hoogte blijven van
ontwikkelingen in digitale technologieën. & Gemiddeld & H &
AI-Impliciet \\
LO5.4.07 & Prioriteit geven aan het identificeren van mogelijkheden om
te leren over digitale technologieën. & Gemiddeld & H & AI-Impliciet \\
LO5.4.08 & Relevante leermogelijkheden identificeren om aan de eigen
digitale competentiebehoeften te voldoen. & Gemiddeld & K &
AI-Impliciet \\
LO5.4.09 & Nauwkeurig de eigen digitale competenties en de eigen
behoeften op dat gebied beoordelen. & Gemiddeld & V & AI-Impliciet \\
LO5.4.10 & Actief deelnemen aan leeractiviteiten om aan de eigen
digitale competentiebehoeften te voldoen. & Gemiddeld & H &
AI-Impliciet \\
LO5.4.11 & Voortdurend digitale technologische ontwikkelingen en de
implicaties daarvan voor de eigen competentiebehoeften en die van
anderen beoordelen. & Geavanceerd & H & AI-Impliciet \\
LO5.4.12 & Zich bezighouden met voortdurende zelfontwikkeling om aan
digitale competentiebehoeften te voldoen. & Geavanceerd & H &
AI-Impliciet \\
LO5.4.13 & Anderen ondersteunen bij het ontwikkelen van zelfvertrouwen
en autonomie in digitale omgevingen. & Geavanceerd & V & AI-Impliciet \\
LO5.4.14 & Beschikbare leermogelijkheden voor digitale competenties
verzamelen voor een specifiek doel. & Geavanceerd & V & AI-Impliciet \\
LO5.4.15 & Zich bezighouden met voortdurende zelfontwikkeling om aan
complexe of gespecialiseerde digitale competentiebehoeften te voldoen. &
Zeer geavanceerd & H & AI-Impliciet \\
LO5.4.16 & Anderen mentoren bij het identificeren en aanpakken van hun
digitale competentiebehoeften. & Zeer geavanceerd & V & AI-Impliciet \\
LO5.4.17 & Leermateriaal ontwerpen om anderen te helpen voldoen aan
complexe of gespecialiseerde digitale competentiebehoeften. & Zeer
geavanceerd & V & AI-Impliciet \\
\end{longtable}

\newpage{}

\section*{Literatuur}\label{literatuur}
\addcontentsline{toc}{section}{Literatuur}

\phantomsection\label{refs}
\begin{CSLReferences}{1}{0}
\bibitem[\citeproctext]{ref-Abendroth-Dias2025}
Abendroth-Dias, K., Arias-Cabarcos, P., Bacco, F. M., Bassani, E.,
Bertoletti, A., et al. (2025). \emph{Generative {AI} Outlook Report -
Exploring the Intersection of Technology, Society and Policy} (E.
Navajas-Cawood, M. Vespe, A. Kotzev, \& R. Van Bavel, Red.).
Publications Office of the European Union.
\url{https://publications.jrc.ec.europa.eu/repository/handle/JRC142598}

\bibitem[\citeproctext]{ref-Boerkamp2024}
Boerkamp, L. G. P., Deursen, A. J. A. M. van, Laar, E. van, Zeeuw, A.
van der, \& Graaf, S. van der. (2024). Exploring Barriers to and
Outcomes of Internet Appropriation Among Households Living in Poverty: A
Systematic Literature Review. \emph{Sage Open}, \emph{14}(1).
\url{https://doi.org/10.1177/21582440241233047}

\bibitem[\citeproctext]{ref-Brundle2025}
Brundle, C., Johansson, J. F., Best, K., Clegg, A., Forster, A.,
Atkinson, T., Foster, M., Humphrey, S., Iliff, A., Inglis, J., Walker,
C., \& Graham, L. (2025). Development of methods to identify digitally
excluded older people, and tailoring of interventions to meet their
digital needs: a protocol for a mixed-methods study (the {INCLUDE}
study). \emph{BMJ open}, \emph{15}(9), e102723.
\url{https://doi.org/10.1136/bmjopen-2025-102723}

\bibitem[\citeproctext]{ref-Centeno2025}
Centeno, C., \& Cosgrove, J. (2025). \emph{Ten Years of DigComp: A
Framework more essential than ever}. Publications Office of the European
Union.
\url{https://publications.jrc.ec.europa.eu/repository/handle/JRC143430}

\bibitem[\citeproctext]{ref-Centeno2024a}
Centeno, C., Cosgrove, J., Cachia, R., Mora, T., et al. (2024a).
\emph{European Digital Skills Certificate ({EDSC}) Feasibility Study}.
Publications Office of the European Union.
\url{https://publications.jrc.ec.europa.eu/repository/handle/JRC138344}

\bibitem[\citeproctext]{ref-Centeno2024b}
Centeno, C., Cosgrove, J., Cachia, R., Mora, T., et al. (2024b).
\emph{European Digital Skills Certificate ({EDSC}) Feasibility Study:
Annexes to the Repoort}. Publications Office of the European Union.
\url{https://publications.jrc.ec.europa.eu/repository/handle/JRC138344}

\bibitem[\citeproctext]{ref-Draghi2024}
Draghi, M. (2024). \emph{The Future of European Competitiveness: A
Competitiveness Strategy for Europe}. European Commission.
\url{https://commission.europa.eu/document/97e481fd-2dc3-412d-be4c-f152a8232961_en}

\bibitem[\citeproctext]{ref-Duckworth2025}
Duckworth, D., \& Fraillon, J. (2025). Computational thinking framework.
In J. Fraillon \& M. Rožman (Red.), \emph{IEA International Computer and
Information Literacy Study 2023: Assessment Framework} (pp. 35--43).
IEA. \url{https://link.springer.com/book/10.1007/978-3-031-61194-0}

\bibitem[\citeproctext]{ref-Commission2012}
European Commission. (2012). \emph{Charter of Fundamental Rights of the
European Union}.
\url{https://eur-lex.europa.eu/eli/treaty/char_2012/oj/eng}

\bibitem[\citeproctext]{ref-Commission2018}
European Commission. (2018). \emph{Council Recommendation of 22 May 2018
on key competences for lifelong learning}.
\url{https://eur-lex.europa.eu/legal-content/EN/TXT/?uri=uriserv:OJ.C_.2018.189.01.0001.01.ENG}

\bibitem[\citeproctext]{ref-Commission2023b}
European Commission. (2023a). \emph{Commission Staff Working Document...
Proposal for a Council Recommendation on improving the provision of
digital skills in education and training}.
\url{https://eur-lex.europa.eu/legal-content/EN/TXT/?uri=CELEX:52023SC0205}

\bibitem[\citeproctext]{ref-Commission2023a}
European Commission. (2023b). \emph{European Declaration on Digital
Rights and Principles for the Digital Decade}.
\url{https://eur-lex.europa.eu/legal-content/EN/TXT/?uri=OJ:JOC_2023_023_R_0001}

\bibitem[\citeproctext]{ref-Commission2025b}
European Commission. (2025a). \emph{Communication on the Union of
Skills}.
\url{https://eur-lex.europa.eu/legal-content/EN/TXT/PDF/?uri=CELEX\%3A52025DC0090}

\bibitem[\citeproctext]{ref-Commission2025a}
European Commission. (2025b). \emph{State of the Digital Decade 2025:
Keep building the {EU's} sovereignty and digital future}.
\url{https://digital-strategy.ec.europa.eu/en/library/state-digital-decade-2025-report}

\bibitem[\citeproctext]{ref-Eurostat2025}
Eurostat. (2025). \emph{Statistics explained: {ICT} specialists in
employment}.
\url{https://ec.europa.eu/eurostat/statistics-explained/index.php?oldid=675865}

\bibitem[\citeproctext]{ref-FariasGaytan2023}
Farias-Gaytan, S., Aguaded, I., \& Ramirez-Montoya, M. (2023). Digital
transformation and digital literacy in the context of complexity within
higher education institutions: a systematic literature review.
\emph{Humanities and Social Sciences Communications}, \emph{10}(1).
\url{https://doi.org/10.1057/s41599-023-01875-9}

\bibitem[\citeproctext]{ref-Kluzer2020}
Kluzer, S., Centeno, C., \& O'Keeffe, W. (2020). \emph{DigComp at Work}.
Publications Office of the European Union.
\url{https://publications.jrc.ec.europa.eu/repository/handle/JRC120376}

\bibitem[\citeproctext]{ref-Lewandowsky2020}
Lewandowsky, S., Smillie, L., Garcia, D., et al. (2020).
\emph{Technology and democracy: Understanding the influence of online
technologies on political behaviour and decision-making}. Publications
Office of the European Union.
\url{https://publications.jrc.ec.europa.eu/repository/handle/JRC122023}

\bibitem[\citeproctext]{ref-MorteNadal2025}
Morte-Nadal, T., \& Esteban-Navarro, M. Á. (2025). Recommendations for
digital inclusion in the use of European digital public services.
\emph{Humanit Soc Sci Commun}, \emph{12}, 273.
\url{https://doi.org/10.1057/s41599-025-04576-7}

\bibitem[\citeproctext]{ref-OnesiOzigagun2024}
Onesi-Ozigagun, N. O., Ololade, N. Y. J., Eyo-Udo, N. N. L., \&
Ogundipe, N. D. O. (2024). Revolutionizing education through {AI}: A
comprehensive review of enhancing learning experiences.
\emph{International Journal of Applied Research in Social Sciences},
\emph{6}(4), 589--607. \url{https://doi.org/10.51594/ijarss.v6i4.1011}

\bibitem[\citeproctext]{ref-Pakhnenko2025}
Pakhnenko, O., Yarovenko, H., Semenog, A., Mordan, Y., \& Tarasenko, O.
(2025). Uncovering patterns of digital transformation of European
economies using self-organising maps. \emph{Problems and Perspectives in
Management}, \emph{23}(3), 581--596.
\url{https://doi.org/10.21511/ppm.23(3).2025.42}

\bibitem[\citeproctext]{ref-Pouliakas2025}
Pouliakas, K., Santangelo, G., \& Dupire, P. (2025). Are artificial
intelligence skills a reward or a gamble? Deconstructing the {AI} wage
premium in Europe. \emph{Eurasian Bus Rev}.
\url{https://doi.org/10.1007/s40821-025-00302-0}

\bibitem[\citeproctext]{ref-Stalmach2023}
Stalmach, A., D'Elia, P., Di Sano, S., \& Casale, G. (2023). Digital
Learning and Self-Regulation in Students with Special Educational Needs:
A Systematic Review of Current Research and Future Directions.
\emph{Education Sciences}, \emph{13}(10), 1051.
\url{https://doi.org/10.3390/educsci13101051}

\bibitem[\citeproctext]{ref-Tarca2024}
Țarcă, V., Luca, F.-A., \& Țarcă, E. (2024). The Digital Edge: Skills
That Matter in the European Labour Market after {COVID-19}.
\emph{Economies}, \emph{12}(10), 273.
\url{https://doi.org/10.3390/economies12100273}

\bibitem[\citeproctext]{ref-Vuorikari2022b}
Vuorikari, R., Jerzak, N., Karpinski, Z., et al. (2022). \emph{Measuring
Digital Skills across the {EU}: Digital Skills Indicator 2.0}.
Publications Office of the European Union.
\url{https://publications.jrc.ec.europa.eu/repository/handle/JRC130341}

\end{CSLReferences}




\end{document}
